\documentclass[12pt,a4paper]{report}
\usepackage[utf8]{inputenc}
\usepackage[ngerman]{babel}
\usepackage{geometry}
\geometry{a4paper, margin=2.5cm}
\usepackage{setspace}
\onehalfspacing
\usepackage{graphicx}
\usepackage{booktabs}
\usepackage{array}
% Linksbündige Spalte mit fester Breite (fuer Tabellen mit viel Text)
\newcolumntype{L}[1]{>{\raggedright\arraybackslash}p{#1}}
\usepackage{amsmath}
\usepackage{listings}
\usepackage{xcolor}
\usepackage{hyperref}
\usepackage{caption}
\usepackage{csquotes}
\usepackage[backend=biber, style=numeric, sorting=none]{biblatex}
\usepackage{titlesec}
\usepackage{tikz}
\usetikzlibrary{positioning,fit,calc,arrows.meta,shapes.geometric,backgrounds}
\usepackage{float}
\hypersetup{hidelinks}
\usepackage{subcaption}
\usepackage{wrapfig}
\usetikzlibrary{decorations.pathreplacing}
% ---- Kapitelüberschriften kompakt ----
\titleformat{\chapter}[hang]
  {\normalfont\LARGE\bfseries}{\thechapter\quad}{0pt}{}
\titlespacing*{\chapter}{0pt}{0pt}{20pt}
% ---- Literatur und Quellen
\addbibresource{Literatur.bib}

% ---- Gliederungstiefe: dritte Ebene nummerieren, Inhaltsverzeichnis bis Ebene zwei ----
\setcounter{secnumdepth}{2}
\setcounter{tocdepth}{2}

% ---- Code-Listing-Style (fuer Python/JS-Ausschnitte) ----
\lstdefinestyle{code}{
    basicstyle=\ttfamily\small,
    breaklines=true,
    frame=single,
    numbers=left,
    numberstyle=\tiny,
    keywordstyle=\color{blue},
    commentstyle=\color{gray},
    stringstyle=\color{teal},
    backgroundcolor=\color{gray!5}
}
\lstset{style=code}
\lstdefinelanguage{JavaScript}{
  keywords={const,let,var,function,return,if,else,for,of,in,new,this,typeof,try,catch,=>},
  keywordstyle=\color{blue}\bfseries,
  ndkeywords={true,false,null,undefined},
  comment=[l]{//}, morecomment=[s]{/*}{*/},
  morestring=[b]', morestring=[b]",
  sensitive=true
}
\renewcommand{\lstlistingname}{Quelltext}
\emergencystretch=2em

\begin{document}

% ==================== TITELSEITE ====================


\begin{titlepage}
    \begin{singlespace}
    \centering
    \vspace*{1.2cm}

    {\Large\bfseries Ostbayerische Technische Hochschule Regensburg\par}
    \vspace{0.3cm}
    {\large Fakultät Informatik und Mathematik\par}

    \vspace{2cm}

    {\LARGE\bfseries
    Multi-Layer Visualisierung von Tiersichtungsdaten
    und Simulation von Tierbewegungen
    in einer Online-Kartenplattform\par}

    \vspace{1.8cm}

    {\Large Bachelorarbeit\par}

    \vspace{0.8cm}

    vorgelegt von\par
    \vspace{0.3cm}
    {\Large\bfseries Yassir Sbai\par}
    \vspace{0.2cm}
    Matrikelnummer: 3341177

    \vspace{1.8cm}

    Betreuer und erster Prüfer:
    Prof. Dr. Jan Dünnweber\par
    \vspace{0.4cm}
    Zweiter Prüfer:
    Prof. Dr. Georgios Raptis\par

    \vfill

    Regensburg, \today

    \end{singlespace}
\end{titlepage}



% ==================== INHALTSVERZEICHNIS ====================

\tableofcontents
\clearpage

\phantomsection
\addcontentsline{toc}{chapter}{\listfigurename}
\listoffigures
\clearpage

\phantomsection
\addcontentsline{toc}{chapter}{\listtablename}
\listoftables
\clearpage


% ====================================================================
% ====================================================================
\chapter{Einleitung}
\label{chap:einleitung}

\section{Motivation und Problemstellung}
\label{sec:motivation}

Wer frei lebende Wildtiere erforscht oder schützt, braucht heute digitale Werkzeuge, die
Bewegungsmuster in Raum und Zeit sichtbar machen. Bei großen Arten wie Bären, bei denen es
zu Konflikten mit Menschen kommen kann, fallen dabei sehr unterschiedliche Daten an:
GPS-Daten besenderter Tiere, automatisch erfasste Kamerabilder und Pfade aus
Bewegungsmodellen. Damit diese Daten sinnvoll ausgewertet werden können, müssen sie nicht nur gespeichert, sondern auch verständlich dargestellt werden - im Regelfall auf einer
interaktiven Karte, auf der räumliche Zusammenhänge unmittelbar zu sehen sind.

Die beiden Datensätze, die in dieser Arbeit untersucht werden, sind nicht nur unterschiedlich formatiert, sie sind von
grundlegend verschiedener Art. Aus einer Langzeitstudie in den rumänischen Karpaten liegen
\emph{Trajektorien} vor: rund 60.000 GPS-Ortungen von 13 besenderten Europäischen
Braunbären (\emph{Ursus arctos}) über vier Jahre. 
Jede GPS-Ortung ist einem der 13 Bären
zugeordnet und mit einem Zeitstempel versehen.


Aus der japanischen Region Fukushima liegen dagegen \emph{Detektionen} vor: zeitgestempelte
Fotoauslösungen fest installierter Wildtierkameras. Welches einzelne Tier ausgelöst hat, ist
nicht bekannt, und damit auch nicht, ob zwei Aufnahmen zum selben Tier gehören. Dieser
Unterschied ist für die Arbeit zentral, weil er bestimmt, welche Auswertungen möglich sind:
Pfade lassen sich nur auf der einen Seite vergleichen, Tagesrhythmen dagegen auf beiden.

An der Ostbayerischen Technischen Hochschule (OTH) Regensburg entstanden in einer
Kooperation mit der Universität Aizu in Japan drei eigenständige Teilsysteme. Das erste ist
eine Kartenanwendung mit zwei getrennten Seiten: eine für die Gerätestandorte und die dort
erfassten Kameraaufnahmen in der Region Aizu (Präfektur Fukushima, Honshū), eine für die
Bewegungspfade der Karpatenbären \cite{LorenzBA}. Das zweite ist ein Modell, das
Bärenbewegungen auf Grundlage der realen GPS-Daten aus den Karpaten simuliert
\cite{BreuBA}. 
Das dritte Teilsystem verarbeitet die Bilder, die von den Wildtierkameras in Japan an ein gemeinsam genutztes E-Mail-Postfach gesendet werden. Dabei werden
Bärenaufnahmen automatisch extrahiert \cite{LelchatBA}. Alle drei entstanden in getrennten studentischen Arbeiten,
und nur das erste hatte eine Benutzeroberfläche.

Daraus ergeben sich drei Einschränkungen.


\begin{enumerate}
    \item Die Daten ließen sich nur getrennt betrachten.
    In der bestehenden Kartenanwendung konnte über die
    Seitenleiste zwischen den beiden Seiten gewechselt
    werden. Jede Seite stellte eine eigene, unabhängige
    Karte auf Basis der JavaScript-Bibliothek Leaflet dar,
    die über den Python-Wrapper Folium eingebunden war.
    Gerätestandorte und Bewegungspfade lagen nie auf
    derselben Karte.

    \item Die bisherige Kartendarstellung war nur bedingt
    für die angestrebte Datenmenge geeignet. Werden alle
    Daten gleichzeitig auf einer Karte dargestellt,
    steigt die Anzahl der Kartenobjekte deutlich an.
    Die bisherige Anwendung verwendete für Bewegungspfade
    und Kreismarker den SVG-Renderer von Leaflet.
    Da diese Objekte einzeln im Dokumentbaum des Browsers
    verwaltet werden, erhöht sich mit ihrer Anzahl der
    Verarbeitungsaufwand. Dadurch kann die Darstellung
    und Bedienung der Karte langsamer werden.

    \item Das Bewegungsmodell nutzte lokale Daten, während
    Kartenanwendung und Bildextraktion auf dieselbe MariaDB
    zugriffen. Dadurch war der Datenzugriff zwischen allen
    Teilsystemen nicht einheitlich.
\end{enumerate}


Das Bewegungsmodell lag zudem nur als eigenständiger Quelltext vor, denn Gegenstand jener
Arbeit war das Modell selbst und nicht seine Einbindung in eine Anwendung. Die erzeugten
Pfade waren damit zwar vorhanden, aber auf keiner Karte sichtbar und nicht gegen die realen
Daten prüfbar. Genau hier setzt diese Arbeit an: Im Ausgangszustand war es nicht möglich, reale Bewegungsmuster, simulierte Bewegungspfade und die Standorte der Wildtierkameras gemeinsam auf einer Karte darzustellen und miteinander zu vergleichen.

\section{Zielsetzung der Arbeit}
\label{sec:zielsetzung}


Ziel dieser Arbeit ist eine einheitliche Kartenanwendung,
die die vorhandenen Daten auf einer gemeinsamen,
mehrschichtigen (Multi-Layer) Hauptkarte darstellt
und das Bewegungsmodell über zusätzliche Experimentseiten
zugänglich macht. Die automatische Bildextraktion bleibt
ein eigenständiger Prozess. Die dabei erfassten
Bärenaufnahmen werden in der gemeinsamen Datenbank
gespeichert und können über die Geräteauswahl in der
Anwendung angezeigt werden. Im Japan-Layer werden
die Kamerastandorte dargestellt, sofern gültige
Koordinaten vorliegen. Für einige der per E-Mail
erfassten Kameras fehlen diese Standortinformationen,
weshalb sie derzeit nicht auf der Karte angezeigt
werden können. Hinzu kommen Werkzeuge zum Vergleich
simulierter Bewegungen mit realen Beobachtungsdaten.
Im Einzelnen sind das die folgenden Teilziele.


\subsection*{Integration und technische Basis}

\begin{itemize}
    
    
    
    \item \textbf{Technologiewechsel der Kartenvisualisierung:}
    Die bisherige Darstellung über Folium und Leaflet
    wird durch MapLibre GL JS ersetzt. Die Bibliothek
    nutzt WebGL zur hardwarebeschleunigten
    Grafikdarstellung im Browser \cite{MatsudaLea2013}
    und unterstützt Vektorkacheln, die geografische
    Geometrien und ihre Attribute statt fertiger
    Kartenbilder enthalten \cite{Rechsteiner2024,MapLibreDocs}.
    Damit sollen auch größere Datenmengen flüssig
    darstellbar werden und zusätzliche
    Darstellungsformen möglich sein, etwa
    Satellitenkarten und dreidimensionales Gelände.
    
    


    \item \textbf{Systemintegration:}
    Die zuvor getrennt dargestellten Daten werden
    auf einer gemeinsamen Hauptkarte mit unabhängig
    steuerbaren Ebenen zusammengeführt. Das
    Bewegungsmodell wird über zusätzliche Experimentseiten
    in die Anwendung eingebunden. Die Ergebnisse der
    eigenständig laufenden Bildextraktion werden
    über die gespeicherten Aufnahmen zugänglich gemacht.
    
    
    \item \textbf{Einheitlicher Datenzugriff:}
    Die vorhandene MariaDB wird weiterhin zur
    Verwaltung der Kamerastandorte und Bärenaufnahmen
    aus Japan sowie der GPS-Ortungen aus Rumänien
    verwendet. Die ursprünglichen Karpaten-GPS-Daten
    liegen zusätzlich als lokale CSV-Datei vor.
    Diese dient als Ersatzdatenquelle, falls die
    Datenbank nicht erreichbar ist.
    
    Die Geschwindigkeitskennwerte des übernommenen
    Bewegungsmodells werden weiterhin aus der lokalen
    SQLite-Datenbank der Vorarbeit gelesen
    \cite{BreuBA}. Da diese Daten dort bereits
    strukturiert vorliegen und über die vorhandenen
    Zugriffsfunktionen genutzt werden können, ist
    eine zusätzliche Übernahme in MariaDB nicht
    erforderlich. Weitere lokale SQLite-Daten
    dienen als Zwischenspeicher für aufbereitete
    Geländeinformationen.


    
    \item \textbf{Interaktions- und Filterkonzept:}
    Forschende sollen einzelne Bären, Zeiträume,
    Jahresphasen und Darstellungsarten auswählen können.
    Außerdem wird die aus dem Ausgangssystem bekannte
    Funktion wiederhergestellt, eigene Marker und Flächen
    auf der Karte anzulegen und wieder zu entfernen.
    Damit können beispielsweise bestimmte Orte oder
    Untersuchungsbereiche markiert und räumlich mit
    den dargestellten Daten verglichen werden.
    
\end{itemize}

\subsection*{Analyse- und Validierungswerkzeuge}

\begin{itemize}
    
    
    \item \textbf{Bewegungsanimation:}
    Für die Karpaten-Daten macht eine Animationskomponente
    den zeitlichen Verlauf der Bewegung eines ausgewählten
    Bären abspielbar. Dabei werden seine reale
    GPS-Trajektorie und die simulierten Bewegungspfade
    für denselben Tag und dieselben Zeitpunkte synchron
    dargestellt. Abweichungen zwischen Modell und
    Beobachtung werden so direkt auf der Karte sichtbar.
    

    \item \textbf{Übertragung des Bewegungsmodells
    auf die japanischen Kameradaten:}
    Die Schrittauswahl des Modells wird erweitert,
    damit es auch ohne vorgegebene GPS-Trajektorie
    Bewegungswege erzeugen kann. Dazu gehören eine
    überarbeitete Richtungswahl und eine optionale
    Gewichtung des Geländes anhand von OpenStreetMap-Daten.
    Für Japan wird ein fester Parameter für die
    Richtungstreue verwendet.
    
   
    \item \textbf{Statistischer Vergleich der
    Tageszeitverteilungen:}
    Da sich aus einzelnen Kameradetektionen keine
    Pfade rekonstruieren lassen, wird der Vergleich
    auf die Zeit verlagert. Über den
    Überlappungskoeffizienten $\Delta$ nach Ridout
    und Linkie \cite{RidoutLinkie2009} wird
    untersucht, wie stark sich die Uhrzeiten
    virtueller Kameraeintritte und realer
    japanischer Fotoaufnahmen überschneiden.
    Die Unsicherheit wird mit einem
    Bootstrap-Verfahren abgeschätzt \cite{Efron1979}.
    In der durchgeführten Auswertung ergibt sich
    eine teilweise zeitliche Übereinstimmung mit
    einem Überlappungskoeffizienten von
    $\Delta = 0{,}522$.
    
\end{itemize}

\subsection*{Abgrenzung des Leistungsumfangs}

Die fachlichen Modelle werden weder vollständig
neu entworfen noch unverändert übernommen.
Übernommen werden die in der Vorarbeit
\cite{BreuBA} aus den Karpaten-Daten
abgeleiteten Geschwindigkeitswerte.

Die Richtungsentscheidung wird erweitert:
Für die Karpaten-Simulation wird die
Richtungstreue aus aufeinanderfolgenden
realen Schritten geschätzt. Die Japan-Simulation
verwendet dagegen einen festen Richtungsparameter,
geschichtete Richtungskandidaten über den
Vollkreis und optional eine Gewichtung
des Geländes.


Nicht Gegenstand dieser Arbeit sind die
Bildverarbeitung und die automatische
Artbestimmung auf Kamerafallenbildern;
ob ein Bild einen Bären zeigt, wird als
bekannt vorausgesetzt. Ebenso ist die
Anwendung nicht für die Verarbeitung
von Echtzeitdaten ausgelegt; eingehende
Bilder werden in Stapeln verarbeitet.

Die Simulation dient nicht der Vorhersage
der tatsächlichen Bewegung eines einzelnen
Tieres. Sie erzeugt stochastische
Bewegungspfade, deren Schrittlängen auf
empirisch abgeleiteten Geschwindigkeitswerten
beruhen. Die Bewegungsrichtungen werden aus
mehreren möglichen Kandidaten durch eine
gewichtete Zufallsauswahl bestimmt.

In der Karpaten-Simulation basiert diese
Gewichtung auf der aus realen GPS-Daten
geschätzten Richtungstreue. In der
Japan-Simulation werden dagegen ein
festgelegter Richtungsparameter und
optional Geländeeigenschaften berücksichtigt.
Ein mit scikit-learn trainiertes Modell
zur Vorhersage individueller Bewegungswege
kommt dabei nicht zum Einsatz.

Die resultierenden Pfade sind somit
mögliche Realisierungen des Bewegungsmodells
und keine Vorhersagen tatsächlich
zurückgelegter Wege. In den Karpaten werden
reale und simulierte Bewegungen zeitlich
synchron gegenübergestellt. Da für Japan
keine realen GPS-Trajektorien einzelner
Bären vorliegen, erfolgt dort ein
statistischer Vergleich der
Tageszeitverteilungen tatsächlicher
Fotoaufnahmen und virtueller
Kameradetektionen. Dieser Vergleich
bewertet die zeitliche Übereinstimmung,
nicht die Genauigkeit einzelner
simulierter Bewegungspfade.



\section{Aufbau der Arbeit}
\label{sec:aufbau}

Die Arbeit ist wie folgt aufgebaut.

Kapitel~\ref{chap:grundlagen} legt die Grundlagen dar. Abschnitt~\ref{sec:geodaten}
behandelt Geodaten und ihre Darstellung im Web, Abschnitt~\ref{sec:bibliotheken} die
infrage kommenden Kartenbibliotheken und ihre Einbindung in den verwendeten
Anwendungsrahmen. Abschnitt~\ref{sec:fachliche-grundlagen} führt die fachlichen Grundlagen
ein, die für die Modellerweiterung und die statistische Auswertung gebraucht werden.

Kapitel~\ref{chap:analyse} analysiert die bestehenden Systeme, ihre Datenbestände und die
übernommenen Modelle. Besonderes Gewicht liegt auf den festgestellten technischen
Einschränkungen, die mehrere Entwurfsentscheidungen in
Kapitel~\ref{chap:implementierung} begründen. Am Ende des Kapitels stehen daraus
abgeleitete, nummerierte Anforderungen.

Kapitel~\ref{chap:implementierung} beschreibt die
Implementierung: die Kartenvisualisierung,
die Integration der Datenquellen und des
Bewegungsmodells, die Bewegungsanimation,
die Übertragung auf die japanischen Kameradaten
sowie den statistischen Vergleich der
Tageszeitverteilungen.

Kapitel~\ref{chap:diskussion} stellt die
Ergebnisse der Funktionsprüfungen und
Simulationsuntersuchungen vor und ordnet ihre
Aussagekraft sowie ihre Grenzen ein.

Kapitel~\ref{chap:fazit} fasst die Arbeit
zusammen, bewertet die Erreichung der
wesentlichen Ziele und gibt einen Ausblick.

% ====================================================================
\chapter{Grundlagen}
\label{chap:grundlagen}

Dieses Kapitel stellt die Grundlagen bereit, auf denen die spätere Umsetzung aufbaut. Die
Abschnitte~\ref{sec:geodaten} und~\ref{sec:bibliotheken} behandeln die technische Seite:
wie Geodaten beschrieben werden, wie Web-Karten sie darstellen und in welchem Rahmen die
Anwendung läuft. Abschnitt~\ref{sec:fachliche-grundlagen} behandelt die fachliche Seite:
wie Tierbewegung modelliert wird und wie sich aus Kamerafallendaten Aussagen gewinnen
lassen.

% =====================================================================
\section{Geodaten und Web-Kartenvisualisierung}
\label{sec:geodaten}

\subsection{Vektor- und Rasterdaten}


Geodaten sind Daten mit Raumbezug. In
Geoinformationssystemen werden zwei grundlegende
Datenmodelle unterschieden: Vektor- und Rasterdaten.

\textbf{Vektordaten} beschreiben räumliche Objekte
durch ihre Geometrie \cite{GISGrundlagen}.
Für diese Arbeit sind drei Typen relevant:


\begin{itemize}
    \item Ein \emph{Punkt} ist eine einzelne Koordinate, zum Beispiel die GPS-Ortung eines
    besenderten Bären oder der Standort eines Geräts.
    \item Eine \emph{Linie} ist eine Folge von Punkten, zum Beispiel der Bewegungspfad eines
    Bären über einen Tag.
    \item Ein \emph{Polygon} ist eine geschlossene Fläche, zum Beispiel ein
    Untersuchungsgebiet oder eine von Nutzenden eingezeichnete Fläche.
\end{itemize}

\textbf{Rasterdaten} bestehen aus einem Gitter von Zellen mit je einem Wert.
Satellitenbilder gehören dazu, ebenso Höhenmodelle und die Kartenbilder, aus denen
Web-Karten zusammengesetzt sind.

Für die Darstellung ist der Unterschied wichtig. Ein Rasterbild wird gezeichnet, wie es
ist. Vektordaten dagegen enthalten keine Angaben zum Aussehen: Farbe, Größe und Form
werden erst beim Zeichnen festgelegt. Deshalb lassen sich dieselben Daten je nach Frage
unterschiedlich darstellen.

\subsection{Koordinatenbezugssysteme und GeoJSON}
\label{subsec:crs}

Eine Koordinate allein ist mehrdeutig. Erst mit dem zugehörigen Bezugssystem bezeichnet
sie einen bestimmten Ort \cite{GISGrundlagen}. Ein solches Koordinatenreferenzsystem legt
fest, wie die Erdfigur angenähert wird, wo der Ursprung liegt und in welcher Einheit und
Reihenfolge die Achsen gezählt werden. Die gebräuchlichen Systeme stehen im \emph{EPSG
Geodetic Parameter Dataset}, einem Katalog der International Association of Oil \& Gas
Producers. Jedes System hat dort eine eindeutige Nummer \cite{EPSG}. In dieser Arbeit
treffen zwei Systeme aufeinander.

Das \textbf{World Geodetic System 1984} (WGS~84, EPSG:4326) ist ein geographisches
Koordinatenreferenzsystem. Seine Achsen sind die geodätische Breite nach Norden und die
geodätische Länge nach Osten, beide in Grad. In diesem System arbeiten GPS-Empfänger, und
GeoJSON schreibt es vor. Alle Daten, die die Karte erreichen, liegen darin vor.

Die Rohdaten der Karpaten-Studie liegen dagegen in \textbf{EPSG:3844} vor, amtlich
\enquote{Pulkovo 1942(58) / Stereo70}. Das ist die rumänische Landesprojektion. Sie bildet
die Erdoberfläche über eine schiefachsige stereographische Projektion auf eine Ebene ab,
mit Ursprungsbreite 46° und Mittelmeridian 25°, und gibt Positionen in Metern an.

Der Unterschied ist grundsätzlich: WGS~84 beschreibt Positionen auf einer gekrümmten
Fläche in Winkeln, EPSG:3844 bildet sie auf eine Ebene ab. Jede solche Abbildung verzerrt.
Eine Projektion kann nur eine Eigenschaft erhalten - Winkel, Flächen oder Entfernungen --
aber nie alle zugleich. Die stereographische Projektion ist winkeltreu und eignet sich
deshalb für Gebiete, die in alle Richtungen etwa gleich weit reichen \cite{Snyder1987}.

Daraus folgt, warum umgerechnet werden muss. Eine Landesprojektion ist auf ein Gebiet
zugeschnitten: Innerhalb Rumäniens sind die Verzerrungen gering und Entfernungen direkt in
Metern ablesbar. Außerhalb wachsen sie schnell an, und für eine Weltkarte ist das System
unbrauchbar.

Die Koordinaten werden bei der Datenaufbereitung nach WGS~84 transformiert.
Dabei erfolgt die Umrechnung sowohl beim Datenbankimport als auch beim Einlesen
der lokalen CSV-Dateien mit \texttt{pyproj}.

Dabei ist ein Detail zu beachten. Das EPSG-Register führt die Achsen von EPSG:4326 in der
Reihenfolge Breite, dann Länge. GeoJSON schreibt die umgekehrte Reihenfolge vor, wie der
folgende Abschnitt zeigt. Die Transformation wird deshalb mit dem Parameter
\texttt{always\_xy} angelegt. Er erzwingt die von GeoJSON verlangte Reihenfolge Länge vor
Breite.

Für die Kartenübergabe werden die Koordinaten einheitlich in WGS~84 und in der
von GeoJSON vorgegebenen Reihenfolge Länge, Breite bereitgestellt.

\paragraph{GeoJSON als Austauschformat.}
\label{subsec:geojson}

Für den Austausch von Vektordaten zwischen Server und Browser wird in dieser Arbeit
durchgängig GeoJSON verwendet. Die Internet Engineering Task Force (IETF) hat es als
RFC~7946 standardisiert \cite{RFC7946}. Es baut auf JSON (JavaScript Object Notation) auf,
einem textbasierten Format aus verschachtelten Schlüssel-Wert-Paaren.

Der Aufbau ist einfach. Ein \texttt{Feature} besteht aus einer Geometrie und einem
Eigenschaftsobjekt (\texttt{properties}) mit beliebigen Schlüssel-Wert-Paaren. Mehrere
Features bilden zusammen eine \texttt{Feature\-Collection}.
Listing~\ref{lst:geojson-grundlagen} zeigt ein Punkt-Feature.

\begin{lstlisting}[caption={Aufbau eines GeoJSON-Punkt-Features},
                   label={lst:geojson-grundlagen}]
{
  "type": "Feature",
  "properties": { "bear_id": "10", "timestamp": "2013-06-14T08:00:00" },
  "geometry": { "type": "Point", "coordinates": [25.041, 45.663] }
}
\end{lstlisting}

Das Format passt aus zwei Gründen zu dieser Arbeit.
Erstens können sowohl das Backend (Python) als auch
das Frontend (JavaScript) GeoJSON direkt verarbeiten:
Python erzeugt die Daten als verschachteltes Dictionary,
JavaScript verarbeitet die übergebene JSON-Repräsentation.
 Zweitens nimmt das freie \texttt{properties}-Objekt fachliche Angaben direkt am
Objekt auf, etwa die Bären-Kennung oder den Zeitstempel. Die Kartenbibliothek kann die
Darstellung später von diesen Angaben abhängig machen.

Zwei Festlegungen des Standards sind praktisch wichtig. Erstens schreibt RFC~7946 WGS~84
als Bezugssystem vor. Zweitens legt er die Reihenfolge der Koordinaten fest: Die ersten
beiden Elemente sind \enquote{longitude and latitude [\ldots] precisely in that order}
\cite{RFC7946}, also Länge vor Breite.

Das ist umgekehrt zu der Reihenfolge, in der man Koordinaten gewöhnlich nennt, und
widerspricht außerdem dem EPSG-Register. Wer die Achsenreihenfolge bei der Umrechnung verwechselt, kann daher Punkte
an einer falschen Position darstellen.

\subsection{Kacheln und Höhenmodelle}
\label{subsec:kacheln}

Eine Web-Karte lädt ihren Hintergrund nicht als ein großes Bild, sondern in quadratischen
Ausschnitten, den \emph{Kacheln}. Viele Webkartendienste verwenden dafür ein gemeinsames Kachelschema.

Auf Zoomstufe~0 zeigt eine einzige Kachel die ganze Welt. Mit jeder weiteren Zoomstufe
wird jede Kachel in vier Teile geteilt, auf Stufe~$z$ gibt es also $4^{z}$ Kacheln. Jede
Kachel hat die Adresse $(z, x, y)$, die direkt in eine URL eingesetzt wird. Der Browser
lädt nur die Kacheln, die gerade sichtbar sind, und holt beim Verschieben die neuen nach.

Das Raster braucht eine Projektion, die sich gleichmäßig in Quadrate teilen lässt. Üblich
ist die Web-Mercator-Projektion (EPSG:3857). Die Anwendung selbst rechnet nicht in diesem
System; die Kartenbibliothek erledigt das intern.


Daraus folgt eine Einschränkung, die leicht übersehen wird:
Kacheldienste stellen Kacheln in der Regel nur bis zu
einer festgelegten maximalen Zoomstufe bereit.
Zoomt man weiter hinein, muss der Browser die
vorhandenen Kacheln vergrößern. Kennt er diese
Obergrenze nicht, fordert er Kacheln an, die es
nicht gibt.


\textbf{Rasterkacheln} sind fertige Bilder, meist PNG oder JPEG. Der Server legt das
Aussehen fest, der Browser zeichnet das Bild nur noch. Satellitenbilder werden in Webkarten üblicherweise als Rasterkacheln bereitgestellt.

\textbf{Vektorkacheln} enthalten statt eines Bildes die Geometrien selbst, in einem binären
Format nach der Mapbox Vector Tile Specification \cite{MVTSpec}. Hier bestimmt erst der
Browser das Aussehen. Farben, Linienstärken und Beschriftungen lassen sich deshalb ändern,
ohne neue Daten zu laden - dieselben Kacheln ergeben eine helle und eine dunkle Karte.

\paragraph{Höhenmodelle.}
\label{subsec:hoehenmodelle}

Für die dreidimensionale Darstellung des Geländes braucht die Karte ein digitales
Höhenmodell (Digital Elevation Model, DEM). Es kommt ebenfalls als Rasterkachel, kodiert
aber keine Farbe, sondern eine Höhe: Die drei Farbkanäle eines Bildpunktes werden gemeinsam
als eine Zahl gelesen, die die Geländehöhe in Metern angibt.

Daraus baut der Browser ein Höhengitter und hebt die Kartengeometrie entsprechend an. Aus
denselben Daten lässt sich eine Schummerung berechnen, die das Relief durch simulierten
Lichteinfall auch in der Draufsicht sichtbar macht. Beide Darstellungsformen dieser Arbeit
beruhen damit auf derselben Datengrundlage.

% =====================================================================
\section{Web-Kartenbibliotheken und ihre Einbindung}
\label{sec:bibliotheken}

Die eigentliche Darstellung der Karte im Browser übernimmt eine JavaScript-Bibliothek.
Dieser Abschnitt stellt die beiden Bibliotheken gegenüber, die für diese Arbeit relevant
sind: Leaflet, auf dem das Ausgangssystem beruhte, und MapLibre GL JS, auf das es
umgestellt wurde. Anschließend wird der Anwendungsrahmen behandelt, in den beide
eingebettet werden.



\subsection{Leaflet und Folium}
\label{subsec:leaflet}

Leaflet ist eine verbreitete, quelloffene JavaScript-Bibliothek für interaktive Karten. Sie
verwaltet die Kacheln des Kartenhintergrunds, reagiert auf Zoomen und Verschieben und
stellt Vektorobjekte wie Marker, Linien und Flächen bereit.

Für das Zeichnen dieser Vektorobjekte verwendet Leaflet standardmäßig einen SVG-Renderer;
über die Option \texttt{preferCanvas} lässt sich alternativ ein Canvas-Renderer wählen
\cite{LeafletDocs}. Im SVG-Fall erhält jedes Objekt ein eigenes Element im Document Object
Model des Browsers. Der Browser muss diese Elemente verwalten und bei jeder Kartenbewegung
neu positionieren; der Aufwand wächst somit mit der Zahl der Objekte. Der Canvas-Renderer
vermeidet die Einzelelemente, zeichnet aber ebenfalls auf der CPU.

Leaflet ist eine reine JavaScript-Bibliothek. Damit sie aus einer Python-Anwendung heraus
nutzbar wird, existiert der Wrapper Folium: Er nimmt Python-Objekte entgegen und erzeugt
daraus eine vollständige HTML-Seite mit eingebettetem Leaflet-Code.

Die Einbindung dieser Seite in eine Streamlit-Anwendung erfolgt über das Paket
\texttt{streamlit-folium}. Es bietet zwei Funktionen an, deren Unterschied für das
Verständnis des Ausgangssystems wesentlich ist \cite{StreamlitFolium}:

\begin{itemize}
    \item \texttt{st\_folium()} wird als \enquote{bi-directional Component} beschrieben.
    Die Funktion liefert ein Dictionary zurück, über das Benutzerinteraktionen - etwa
    gezeichnete Objekte oder angeklickte Elemente - an Python zurückgemeldet werden.
    \item \texttt{folium\_static()} stellt die Karte über deren HTML-Repräsentation dar und
    ist laut Dokumentation eine echte Teilmenge der Funktionalität von
    \texttt{st\_folium()}. Es gibt keinen Rückgabewert; Benutzereingaben erreichen die
    Anwendung nicht.
\end{itemize}

Eine mit \texttt{folium\_static()} eingebundene Karte kann folglich nichts entgegennehmen,
was der Nutzer auf ihr tut. Dieser Unterschied erklärt die in
Kapitel~\ref{chap:einleitung} beschriebene Ungleichheit der beiden Ausgangskarten und wird
in Kapitel~\ref{chap:analyse} anhand des Quellcodes belegt.

\subsection{MapLibre GL JS und seine Eigenschaften}
\label{subsec:maplibre}

Mapbox GL JS geht einen anderen Weg. Die Bibliothek verwaltet Vektorobjekte nicht als
Elemente im Dokumentbaum, sondern übergibt sie als Geometriedaten an die Grafikkarte und
lässt sie dort über die Browser-Schnittstelle WebGL zeichnen. Damit wurden Vektorkacheln
im Browser praktisch nutzbar.

Im Dezember 2020 stellte Mapbox die Bibliothek auf eine nicht quelloffene Lizenz um. Aus
dem letzten frei lizenzierten Stand entstand MapLibre GL JS als quelloffener Fork, zunächst
als direkter Ersatz gedacht. MapLibre GL JS steht unter der 3-Clause-BSD-Lizenz
\cite{MapLibreRepo}. Die eigene Dokumentation beschreibt es als \enquote{a TypeScript
library that uses WebGL to render interactive maps from vector tiles in a browser}
\cite{MapLibreDocs}.


Bei der Lizenzierung ist zwischen der
Kartenbibliothek und den eingebundenen
Kartendiensten zu unterscheiden. Während die
Lizenz von MapLibre GL JS die Verwendung der
Bibliothek regelt, gelten für die bereitgestellten
Karten- und Höhendaten die Bedingungen der
jeweiligen Anbieter.

Für die Straßenkartendarstellung wird OpenFreeMap genutzt.
Der Anbieter stellt auf Grundlage von OpenStreetMap-Daten
verschiedene Kartenstile bereit. In der Anwendung werden
die Stile Liberty, Bright, Positron und Dark verwendet.
Der öffentliche Dienst ist kostenlos und erfordert weder
eine Registrierung noch einen API-Schlüssel.
Die vorgeschriebenen Quellenangaben zu den Kartendaten
werden über die Kartenansicht angezeigt
\cite{OpenFreeMap}.

Für die Satellitenansicht wird Esri World Imagery
verwendet. Der Dienst stellt Satelliten- und Luftbilder
bereit und unterliegt den Nutzungsbedingungen von Esri.
Die Quellenangaben zu Esri und den beteiligten
Bilddatenanbietern werden in der Kartenansicht angezeigt
\cite{EsriTerms}.

Für die optionale dreidimensionale Geländedarstellung werden
Höhenkacheln des Projekts Mapterhorn eingebunden. Diese
werden aus verschiedenen offenen Geodatenbeständen
zusammengestellt. Die zugehörigen Datenquellen und
Quellenangaben dokumentiert der Anbieter gesondert
\cite{MapterhornAttribution}.

Abbildung~\ref{fig:grundkartenwahl} zeigt
die Grundkartenauswahl der entwickelten
MapLibre-Anwendung. Neben der dargestellten
Satellitenansicht können weitere Kartenstile
ausgewählt werden.

\begin{figure}[htbp]
    \centering
    \includegraphics[width=0.95\linewidth]
    {grundkarten-auswahl.png}
    \caption{Grundkartenauswahl der
    MapLibre-basierten Hauptkarte.
    Dargestellt ist die Satellitenansicht
    von Esri World Imagery mit geöffneter
    Auswahlliste der verfügbaren Kartenstile.
    Eigene Darstellung.}
    \label{fig:grundkartenwahl}
\end{figure}


\paragraph{Datenquellen und Ebenen.}

MapLibre GL JS trennt Daten und Darstellung. Das ist sein zentrales Gestaltungsprinzip.

Eine \emph{Datenquelle} (\texttt{source}) verweist auf einen Datenbestand, etwa eine
GeoJSON-FeatureCollection oder einen Vektorkachel-Dienst. Sie legt fest, \emph{was}
dargestellt wird, nicht wie es aussieht.

Eine \emph{Ebene} (\texttt{layer}) greift auf eine Datenquelle zu und legt fest, \emph{wie}
deren Inhalt gezeichnet wird: als Kreis, als Linie oder als gefüllte Fläche, jeweils mit
Farbe und Größe. Mehrere Ebenen können dieselbe Datenquelle unterschiedlich darstellen.

Wichtig ist, dass das Aussehen von den Eigenschaften der Daten abhängen darf. Eine Ebene
kann die Farbe eines Punktes aus dessen \texttt{properties} ableiten, statt sie fest
vorzugeben. Eine Unterscheidung wie die zwischen gemessenen und simulierten Punkten wird so
einmal in den Daten hinterlegt und von der Karte automatisch umgesetzt.

Die gesamte Kartenkonfiguration steht in einem Stildokument, dessen Aufbau die MapLibre
Style Spec festlegt \cite{MapLibreStyleSpec}. Eine Änderung daran wirkt sofort, ohne dass
Daten neu geladen werden.

MapLibre GL JS unterstützt außerdem Höhenmodelle als Datenquelle und eine neigbare Kamera.
Damit ist die dreidimensionale Darstellung des Geländes möglich \cite{MapLibreDocs}.

\paragraph{Gegenüberstellung der Bibliotheken.}
\label{subsec:vergleich}

Tabelle~\ref{tab:bibliotheken} stellt die wichtigsten Unterschiede gegenüber. Verglichen
werden Funktionsweise und vorhandene Funktionen; wie sich die Darstellung in der Praxis
schlägt, bewertet Kapitel~\ref{chap:diskussion}.

\begin{table}[htbp]
\centering
\caption{Gegenüberstellung der betrachteten Kartenbibliotheken}
\label{tab:bibliotheken}
\small
\begin{tabular}{p{3.8cm} p{4.6cm} p{4.6cm}}
\toprule
\textbf{Kriterium} & \textbf{Leaflet (über Folium)} & \textbf{MapLibre GL JS} \\
\midrule
Zeichnen von Vektorobjekten
  & SVG (Standard) oder Canvas, jeweils auf der CPU
  & WebGL, auf der Grafikkarte \\
\addlinespace
Verwaltung einzelner Objekte
  & Bei SVG ein Element je Objekt im Dokumentbaum
  & Keine Einzelelemente im Dokumentbaum \\
\addlinespace
Kachelarten
  & Rasterkacheln
  & Raster- und Vektorkacheln \\
\addlinespace
Änderung des Aussehens zur Laufzeit
  & Neuaufbau der betroffenen Objekte
  & Eigenschaft im Stildokument ändern \\
\addlinespace
Dreidimensionale Geländedarstellung
  & Nicht als Standardfunktion enthalten
  & Über Höhenmodell als Datenquelle \\
\addlinespace
Einbindung in die bestehende Anwendung
  & Folium und streamlit-folium
  & Eigene JavaScript-Einbindung \\
\bottomrule
\end{tabular}
\end{table}

MapLibre GL JS war für diese Arbeit als Kartenbibliothek vorgegeben. Seine Trennung
von Datenquellen und Darstellungsebenen passt zur geforderten Multi-Layer-Karte.
Die Einbindung in Streamlit erfolgt über eine eigene JavaScript-Komponente,
die in Kapitel~\ref{chap:implementierung} beschrieben wird.

\subsection{Streamlit als Anwendungsrahmen}
\label{subsec:streamlit}

Die Anwendung, mit der sich diese Arbeit befasst, ist in Python geschrieben und nutzt Streamlit für die Benutzeroberfläche. Streamlit erzeugt aus einem gewöhnlichen Python-Skript eine
Webanwendung: Aufrufe wie \texttt{st.selectbox()} oder \texttt{st.button()} erzeugen
Bedienelemente, deren Werte das Skript unmittelbar weiterverarbeiten kann.

Drei Eigenschaften dieses Rahmens bestimmen den Aufbau der späteren Implementierung und
werden deshalb hier eingeführt.

\paragraph{Das Ausführungsmodell.}


Streamlit verwendet ein skriptbasiertes Ausführungsmodell.
Bei üblichen Interaktionen mit Bedienelementen wird das
Python-Skript erneut von Anfang bis Ende ausgeführt
\cite{StreamlitCaching}. Formulare bilden eine Ausnahme:
Sie können mehrere Eingaben zusammenfassen, sodass die
Verarbeitung erst beim Absenden ausgelöst wird.

Dieses Modell macht den Programmablauf einfach nachvollziehbar, hat aber eine unmittelbare
Folge: Jede aufwendige Berechnung im Skript wird bei jeder Interaktion wiederholt - auch
dann, wenn sie mit der ausgelösten Änderung nichts zu tun hat. Eine Datenbankabfrage für
die eine Datenebene läuft also auch dann erneut, wenn der Nutzer einen Filter einer ganz
anderen Datenebene verstellt.

\paragraph{Zwischenspeicherung.}

Gegen dieses Problem stellt Streamlit zwei Dekoratoren bereit, die sich in ihrem
Einsatzzweck unterscheiden \cite{StreamlitCaching}:

\begin{itemize}
    \item \texttt{st.cache\_data} ist für Daten gedacht, also für serialisierbare
    Rückgabewerte wie Tabellen, Zahlenfelder oder Zeichenketten. Jeder Aufruf erhält eine
    eigene Kopie des Ergebnisses, wodurch ungewollte Veränderungen durch einen Aufrufer
    andere Aufrufer nicht beeinflussen.
    \item \texttt{st.cache\_resource} ist für globale Betriebsmittel gedacht, die sich
    nicht serialisieren lassen - etwa Datenbankverbindungen oder geladene Modelle. Alle
    Durchläufe und alle Sitzungen teilen sich dasselbe Objekt, ohne es zu kopieren. Das
    Objekt muss daher nebenläufigkeitssicher sein, und Änderungen daran wirken sich auf
    alle Nutzer aus.
\end{itemize}

In der Anwendung werden vor allem aufbereitete Daten und berechnete Ergebnisse
zwischengespeichert.

\paragraph{Einbindung eigener Komponenten.}
\label{par:komponenten}

Streamlit bringt kein Kartenwidget mit dem benötigten Funktionsumfang mit. MapLibre GL JS
muss deshalb als eigene Komponente eingebunden werden. Dafür bestehen zwei Wege, die sich
grundlegend unterscheiden \cite{StreamlitComponents}.

Gemeinsam ist beiden, dass die Komponente als eigenständige Webseite in einem
\texttt{iframe} dargestellt wird. Sie läuft damit in einem eigenen Dokument und hat keinen
direkten Zugriff auf die umgebende Streamlit-Seite.

\textbf{Statische Komponenten} werden über \texttt{components.html()} erzeugt. Die Funktion
nimmt eine Zeichenkette mit HTML entgegen und bettet sie in einen \texttt{iframe} ein. Die
Übertragung verläuft nur in eine Richtung: Die Komponente zeigt an, was Python ihr
mitgibt, meldet aber nichts zurück.

\textbf{Bidirektionale Komponenten} werden über \texttt{declare\_component()} erzeugt. Beim
Aufruf übergebene Argumente erreichen die Komponente; ruft deren JavaScript-Code
\texttt{Streamlit.setComponentValue()} auf, gelangt der übergebene Wert zurück nach
Python. Streamlit führt daraufhin das Skript erneut aus, und der Aufruf der Komponente
liefert in diesem Durchlauf den neuen Wert. Da nur ein einzelner Wert übergeben werden
kann, müssen mehrere Angaben in einem Objekt zusammengefasst werden.


% =====================================================================
\section{Bewegungsmodellierung und Kamerafallendaten}
\label{sec:fachliche-grundlagen}

Dieser Abschnitt führt die fachlichen Grundlagen ein, die für die Erweiterung des
Bewegungsmodells und für den statistischen Vergleich gebraucht werden.

\subsection{Bewegungsökologie und Step-Selection-Funktionen}
\label{subsec:bewegungsoekologie}


\paragraph{Trajektorie, Schritt und Drehwinkel.}

In der Bewegungsökologie wird die Bewegung eines Tieres als Folge einzelner Schritte
beschrieben \cite{Thurfjell2014}. Ein \emph{Schritt} verbindet zwei aufeinanderfolgende
Positionen $x_{t-1}$ und $x_t$ im Zeitabstand
$\Delta t$. Zur Beschreibung dieser Schritte
werden insbesondere zwei Größen verwendet
\cite{Thurfjell2014}:

\begin{itemize}
    \item die \emph{Schrittlänge} $l$, also die
    geradlinige Entfernung zwischen den beiden
    Positionen,
    \item der \emph{Drehwinkel} $\theta$, also die
    Richtungsänderung gegenüber dem vorherigen Schritt.
\end{itemize}


Für die Bewegungsmodellierung ist neben der
Schrittlänge und dem Drehwinkel auch der
Zeitabstand zwischen zwei Positionsmessungen
relevant. Bei GPS-Telemetriedaten können diese
Abstände beispielsweise durch Messlücken
unterschiedlich groß sein. In der Karpaten-Animation, mit der sich diese Arbeit befasst, wird die Länge eines simulierten Schritts aus der angenommenen Geschwindigkeit und dem
tatsächlich verstrichenen Zeitabstand zwischen zwei GPS-Positionen berechnet.



\paragraph{Korrelierter Zufallspfad.}

Das einfachste Modell ist der einfache Zufallspfad: Jeder Schritt bekommt eine neue,
gleichverteilt gezogene Richtung. Echte Tiere laufen nicht so. Sie behalten eine Richtung
über mehrere Schritte bei. Diese Eigenschaft heißt Richtungstreue.

Der \emph{korrelierte Zufallspfad} bildet sie ab: Der Drehwinkel wird nicht
gleichverteilt gezogen, sondern aus einer Verteilung, die sich um null häuft. Eine häufig verwendete Kreisverteilung zur Modellierung
von Drehwinkeln in korrelierten Zufallspfaden ist die von-Mises-Verteilung \cite[S.~818, Box~1]{Codling2008}.
Ihre Wahrscheinlichkeitsdichte lautet

\begin{equation}
    p(\theta) \;=\; \frac{1}{2\pi I_{0}(\kappa)}\,
    \exp\!\bigl(\kappa \cos(\theta - \mu)\bigr),
    \label{eq:vonmises}
\end{equation}

mit der mittleren Richtung $\mu$, dem Konzentrationsparameter $\kappa$ und der
modifizierten Besselfunktion erster Art nullter Ordnung $I_{0}$.

$\kappa$ steuert das Verhalten vollständig. Bei $\kappa = 0$ ist die Verteilung über den
Kreis gleichmäßig, das Tier läuft also richtungslos umher. Je größer $\kappa$, desto stärker
konzentriert sich die Verteilung auf die bisherige Richtung und desto gestreckter wird die
Bewegung.

Wichtig für die Anwendung: $\kappa$ lässt sich aus Daten schätzen. Für beobachtete
Drehwinkel $\theta_{1},\dots,\theta_{n}$ ist der mittlere Kosinus

\begin{equation}
    \rho \;=\; \frac{1}{n}\sum_{j=1}^{n} \cos\theta_{j}.
    \label{eq:rho}
\end{equation}

Bei einer um null zentrierten von-Mises-Verteilung gilt für nichtnegative
Konzentrationsparameter $\kappa$ die Beziehung
$A(\kappa)=I_{1}(\kappa)/I_{0}(\kappa)=\rho$. Die Umkehrung wird in
der Praxis näherungsweise berechnet. Negative Werte von $\rho$ werden in
der implementierten Richtungsgewichtung durch eine Bevorzugung von Umkehrungen behandelt.

Anschaulich misst $\rho$, wie geradlinig die Bewegung ist: $+1$ heißt, jeder Schritt setzt
die Richtung des vorherigen genau fort, $0$ heißt richtungslos, negative Werte heißen
überwiegend Hin und Her.

\paragraph{Step-Selection-Funktionen.}

Step-Selection-Funktionen (SSF) verbinden Bewegung mit Umweltpräferenz. Das Prinzip ist
einfach: Jedem beobachteten Schritt werden $k$ \emph{mögliche} Schritte gegenübergestellt.
Sie starten am selben Punkt, führen aber zu zufällig gezogenen Zielen. Aus dem Vergleich
der Umwelteigenschaften am echten Ziel mit denen an den Alternativen lässt sich schätzen,
was bevorzugt und was gemieden wird \cite{Thurfjell2014}.

Die Methode ist in der Bewegungsökologie etabliert und wird unter anderem für Habitatwahl,
Wanderkorridore und Mensch-Wildtier-Konflikte eingesetzt \cite{Thurfjell2014}. Wichtig ist
dabei, wie die Alternativschritte gezogen werden: Das Ziehverfahren beeinflusst die
Schätzung systematisch \cite{Forester2009}.


\paragraph{Einordnung für diese Arbeit.}

Die klassische Step-Selection-Analyse dient dazu, aus
beobachteten und alternativ erzeugten Bewegungsschritten
Habitatpräferenzen statistisch zu schätzen, beispielsweise
mittels bedingter logistischer Regression. Eine solche
Schätzung von Präferenzkoeffizienten wird in dieser Arbeit
nicht durchgeführt. Aus den Karpaten-Daten wird dagegen die
Richtungstreue geschätzt, also eine Eigenschaft der Bewegung
selbst und keine Umweltpräferenz. Die ermittelten Werte
werden als feste Stützwerte im Modell hinterlegt. Die
Geschwindigkeitswerte stammen unverändert aus der Vorarbeit
von Bräu \cite{BreuBA}.

Stattdessen wird das Prinzip der Erzeugung alternativer
Bewegungsschritte für die Simulation aufgegriffen. Aus den
erzeugten Kandidaten wird jeweils ein Schritt zufällig,
aber gewichtet ausgewählt. Dabei werden die Richtungstreue
und in der Japan-Simulation optional die
Geländeeigenschaften berücksichtigt. Durch die wiederholte
Auswahl entstehen simulierte Bewegungspfade.

Diese Pfade sind keine statistischen Schätzungen von
Habitatpräferenzen, sondern Realisierungen eines
Bewegungsmodells, dessen Parameter teilweise aus den
Karpaten-Daten abgeleitet und teilweise festgelegt wurden.

\subsection{Gewichtete Schwerpunktbildung}
\label{subsec:schwerpunkt}

Ein letzter Baustein stammt nicht aus der Bewegungsökologie, sondern aus einem anderen
Anwendungsgebiet.


Ein mögliches Verfahren zur Lokalisierung von
Fußpositionen auf Bodensensormatten wird in
\cite{GeyerMelzerDuennweber2026} beschrieben.
Dabei wird die Sensorfläche in Felder unterteilt,
die entlang ihrer Diagonalen nochmals in
Dreiecke zerlegt werden. Betritt ein Fuß die
Fläche, werden mehrere dieser Dreiecke
unterschiedlich stark aktiviert.

Die Position wird dann nicht als Mittelpunkt
der aktivierten Fläche bestimmt, sondern als
gewichteter Schwerpunkt:


\begin{equation}
    P \;=\; \frac{\sum_{i=1}^{n} w_{i}\, p_{i}}{\sum_{i=1}^{n} w_{i}} .
    \label{eq:schwerpunkt}
\end{equation}

Dabei ist $P$ die gesuchte Position, $p_{i}$ der Mittelpunkt des $i$-ten aktivierten
Dreiecks und $w_{i}$ dessen Aktivierungswert. Die Gewichtung bewirkt, dass die berechnete
Position sich zu den stärker aktivierten Bereichen hin verschiebt, also dorthin, wo der
größere Teil des Fußes aufliegt.

Für die vorliegende Arbeit ist weniger der konkrete Anwendungsfall von Bedeutung als die
zugrunde liegende Idee: Aus einer Menge unterschiedlich stark bewerteter Flächenstücke wird
ein einzelner, nach diesen Bewertungen gewichteter Bezugspunkt bestimmt. Dieselbe Aufgabe
stellt sich bei der Richtungswahl in strukturiertem Gelände, wenn aus einer Menge
unterschiedlich attraktiver Geländezellen ein Ziel abgeleitet werden soll.
Kapitel~\ref{chap:implementierung} beschreibt die Übertragung.

\subsection{Kamerafallen und Aktivitätsmuster}
\label{subsec:kamerafallen}

\paragraph{Detektionen statt Trajektorien.}

Kamerafallen liefern eine ganz andere Datenform als GPS-Telemetrie. Eine Kamera löst aus,
wenn sich vor ihr etwas bewegt, und speichert Bild und Zeitpunkt. Daraus folgen drei
Einschränkungen.


\begin{enumerate}
    \item Das fotografierte Tier lässt sich meist nicht
    als Individuum bestimmen. Zwei Aufnahmen können
    dasselbe Tier zeigen oder zwei verschiedene.

    \item Zwischen zwei Aufnahmen ist kein
    Bewegungszusammenhang bekannt. Aus Detektionen an
    verschiedenen Kameras lässt sich kein Weg
    rekonstruieren, weil weder feststeht, ob es
    dasselbe Tier war, noch welchen Weg es genommen hat.

    \item Es wird nur punktuell beobachtet. Erfasst
    wird, was sich im Bereich der Kamera bewegt;
    über alles dazwischen gibt es keine Information.
\end{enumerate}


Deshalb gibt es in dieser Arbeit keinen räumlichen Pfadvergleich zwischen Simulation und
japanischen Realdaten: Auf der Seite der Realdaten gibt es keine Pfade, mit denen man
vergleichen könnte. Der Vergleich wird stattdessen auf die Zeit verlagert.

\paragraph{Aktivitätsmuster als zirkuläre Daten.}

Auch ohne Individualerkennung können die Auslösezeiten Hinweise auf das
zeitliche Aktivitätsmuster geben. Die Erfassungen hängen jedoch unter anderem
vom Kamerastandort und der Erfassungswahrscheinlichkeit ab \cite{Rowcliffe2014}.

Uhrzeiten sind dabei keine gewöhnlichen Zahlen, sondern \emph{zirkuläre} Daten. Der Tag ist
ein geschlossener Kreis: 23:59~Uhr und 00:01~Uhr liegen zwei Minuten auseinander, als Zahlen
aber fast 24~Stunden. Wer das nicht berücksichtigt, bekommt genau an der Tagesgrenze falsche
Ergebnisse - besonders bei nachtaktiven Arten, bei denen sich viele Beobachtungen um
Mitternacht häufen.

Die Uhrzeit wird deshalb in einen Winkel umgerechnet, wobei 24~Stunden einem vollen Umlauf
von $2\pi$ entsprechen. Aus diesen Winkeln wird dann eine Dichtefunktion geschätzt, die die
Aktivität über den Tag beschreibt. Dafür kann eine periodisch fortgesetzte Kerndichteschätzung verwendet werden
\cite{RidoutLinkie2009}.




\paragraph{Der Überlappungskoeffizient $\Delta$.}

Liegen zwei solche Dichtefunktionen vor - etwa für zwei Arten oder, wie hier, für
Realdaten und Simulation -, stellt sich die Frage, wie ähnlich sie sind. Ridout und
Linkie schlagen dafür den Überlappungskoeffizienten $\Delta$ vor
\cite{RidoutLinkie2009}.

Anschaulich ist $\Delta$ die Fläche, die sich beide Kurven teilen.
Abbildung~\ref{fig:aktivitaetskurven-beispiel} zeigt das an zwei schematischen
Tagesaktivitätsverteilungen.

\begin{figure}[H]
    \centering
    \includegraphics[width=14cm]{abb_delta_schema_ohne_formel.pdf}
    \caption{Schematische Darstellung zweier Tagesaktivitätsverteilungen.
    Die grau markierte Fläche unter der jeweils niedrigeren der beiden Kurven
    entspricht dem Überlappungskoeffizienten $\Delta$.
    Eigene Darstellung mit künstlichen Beispielwerten.}
    \label{fig:aktivitaetskurven-beispiel}
\end{figure}

In formaler Schreibweise:

\begin{equation}
    \Delta \;=\; \int_{0}^{24} \min\bigl(f(t),\, g(t)\bigr)\, \mathrm{d}t .
    \label{eq:delta}
\end{equation}

Der Wert liegt zwischen 0 und 1: $\Delta = 0$ heißt, die Aktivitätszeiten
überschneiden sich nicht, $\Delta = 1$ heißt, beide Verteilungen sind identisch.


Für $\Delta$ stehen verschiedene Schätzverfahren zur Verfügung.
Ihre Eignung hängt unter anderem vom Stichprobenumfang ab.


Das Maß wird für den Vergleich von Tagesrhythmen eingesetzt. Wichtig ist, was es
\emph{nicht} sagt: $\Delta$ misst nur die zeitliche Übereinstimmung. Über das
räumliche Verhalten, über zurückgelegte Strecken oder darüber, ob sich zwei Tiere
wirklich begegnen, sagt der Wert nichts.


\paragraph{Stichprobenumfang und Unsicherheit.}

$\Delta$ wird aus einer begrenzten Zahl von Beobachtungen geschätzt und ist deshalb selbst
unsicher. Zwei Punkte sind dabei wichtig.


Zum einen hängt die Stabilität der geschätzten
Aktivitätskurven vom Stichprobenumfang ab.
In einer Untersuchung werden rund 100~Detektionen
als Richtwert für stabilere Aktivitätskurven genannt
\cite{Cove2018}. Dieser Wert ist keine allgemeine
Mindestgrenze.

Für den Vergleich in dieser Arbeit stehen
49 reale Fotozeitpunkte aus Japan zur Verfügung.
Damit liegt die Stichprobe unter dem genannten
Richtwert. Die Ergebnisse müssen deshalb
mit entsprechender Vorsicht interpretiert werden.


Zum anderen sollte man einen Einzelwert nie ohne Unsicherheitsangabe berichten. Üblich ist das Bootstrap-Verfahren \cite{Efron1979}: Aus den vorhandenen Beobachtungen wird wiederholt eine neue
Stichprobe gleicher Größe mit Zurücklegen gezogen, für jede wird $\Delta$ neu berechnet, und
aus der Verteilung dieser Werte ergibt sich ein Konfidenzintervall. Das Verfahren benötigt keine vorgegebene parametrische Verteilung; die
zeitliche Abhängigkeit von Aufnahmen muss bei der Interpretation jedoch
berücksichtigt werden.

Beide Punkte kommen in Kapitel~\ref{chap:implementierung} bei der Berechnung und in
Kapitel~\ref{chap:diskussion} bei der Einordnung wieder vor.

% ====================================================================

% ====================================================================
% ====================================================================
\chapter{Analyse der bestehenden Systeme und Ableitung der Anforderungen}
\label{chap:analyse}

Dieses Kapitel beschreibt den Zustand der Systeme zu Beginn dieser Arbeit. Es behandelt die vorhandenen Komponenten, die Datenbestände sowie die übernommenen Modelle. Auf der Grundlage der identifizierten Einschränkungen werden anschließend die Anforderungen an das zu entwickelnde System abgeleitet.

% =====================================================================
\section{Ausgangszustand des Systems}
\label{sec:ausgangszustand}

\subsection{Überblick über die Teilsysteme}
\label{subsec:systemuebersicht}

Zu Beginn dieser Arbeit lagen die Ergebnisse von drei Vorarbeiten vor, die unabhängig
voneinander entstanden sind. Nur eine davon war eine Anwendung mit Benutzeroberfläche.
Tabelle~\ref{tab:teilsysteme} nennt die drei Arbeiten,
Abbildung~\ref{fig:ausgangszustand} zeigt, auf welche Daten sie zugreifen.

\begin{table}[htbp]
\centering
\caption{Übersicht der vorgefundenen Vorarbeiten}
\label{tab:teilsysteme}
\small
\begin{tabular}{L{2.9cm} L{5.2cm} L{4.6cm}}
\toprule
\textbf{Vorarbeit} & \textbf{Gegenstand} & \textbf{Datenhaltung} \\
\midrule
Marcel Lorenz \cite{LorenzBA}
  & Kartenanwendung mit zwei Seiten: Gerätestandorte in Japan sowie
    Bewegungspfade der Karpatenbären
  & MariaDB; Bilder als Dateien im Dateisystem, in der Datenbank nur deren Pfade \\
\addlinespace
Fabian Bräu \cite{BreuBA}
  & Modell zur Erzeugung alternativer Bewegungsschritte
  & eigene SQLite-Datei und CSV-Datei \\
\addlinespace
Anthony Lelchat \cite{LelchatBA}
  & Übernahme neuer Kameraaufnahmen aus einem E-Mail-Postfach in die Datenbank
  & dieselbe MariaDB wie die Kartenanwendung \\
\bottomrule
\end{tabular}
\end{table}

\begin{figure}[H]
\centering
\begin{tikzpicture}[
    font=\footnotesize,
    >={Stealth[length=2.2mm,width=1.6mm]},
    every node/.style={align=center, inner sep=2pt},
    box/.style={draw=black!80, fill=white, line width=0.5pt,
                minimum height=10mm, text width=#1},
    box/.default=34mm,
    db/.style={cylinder, shape border rotate=90, aspect=0.16, draw=black!80,
               fill=black!4, line width=0.5pt, minimum width=#1,
               minimum height=12mm, text width=#1},
    db/.default=24mm,
    flow/.style={->, line width=0.6pt, draw=black!85},
    ref/.style={->, densely dashed, line width=0.5pt, draw=black!60},
    lbl/.style={font=\scriptsize, text=black!75, inner sep=1.5pt, fill=white},
    head/.style={font=\small\bfseries, text width=40mm},
]

\def\xA{3.1}
\def\xB{8.3}
\def\xC{12.9}
\def\xDB{5.7}

% --- Spaltenköpfe ---
\node[head] (hA) at (\xA,0)
  {Kartenanwendung\\[-1pt]\normalfont\scriptsize\color{black!60}Vorarbeit Lorenz \cite{LorenzBA}};
\node[head] (hB) at (\xB,0)
  {Bildextraktion\\[-1pt]\normalfont\scriptsize\color{black!60}Vorarbeit Lelchat \cite{LelchatBA}};
\node[head] (hC) at (\xC,0)
  {Bewegungsmodell\\[-1pt]\normalfont\scriptsize\color{black!60}Vorarbeit Bräu \cite{BreuBA}};

% --- Reihe 1 ---
\node[box=40mm, minimum height=12mm] (app) at (\xA,-1.75)
  {\textbf{Streamlit-Anwendung}\\[-1pt]\scriptsize zwei getrennte Kartenseiten:\\[-2pt]\scriptsize Japan und Karpaten};
\node[box=34mm, rounded corners=4pt, fill=black!4] (mail) at (\xB,-1.75)
  {E-Mail-Postfach\\[-1pt]\scriptsize Microsoft Graph API};

% --- Reihe 2 ---
\node[box=40mm] (wdb) at (\xA,-3.85)
  {\ttfamily WildlifeDatabase\\[-1pt]\normalfont\scriptsize Zugriff auf die MariaDB};
\node[box=37mm] (ext) at (\xB,-3.85)
  {\scriptsize\ttfamily WildlifePicture-\\[-2pt]\scriptsize\ttfamily Extractor\\[-1pt]\normalfont\scriptsize Texterkennung, Dauerprozess};
\node[box=37mm] (ssm) at (\xC,-3.85)
  {\scriptsize\ttfamily StepSelectionModel\\[-1pt]\normalfont\scriptsize nur Kandidatenschritte};

% --- Reihe 3 ---
\node[db=30mm] (maria) at (\xDB,-5.95)
  {\textbf{MariaDB}\\[-1pt]\scriptsize zentrale Instanz,\\[-2pt]\scriptsize 11 Tabellen};
\node[db=22mm] (fs)     at (1.15,-5.95)      {Dateisystem\\[-1pt]\scriptsize Bilddateien};
\node[db=22mm] (sqlite) at (\xC-1.15,-5.95)  {SQLite\\[-1pt]\scriptsize\ttfamily database.db};
\node[db=22mm] (csv)    at (\xC+1.15,-5.95)  {CSV\\[-1pt]\scriptsize\ttfamily 1\_bears\_RO};

% --- Verbindungen ---
\draw[flow, <->] (app) -- (wdb) node[lbl, midway] {lesen/schreiben};
\draw[flow, <->] (wdb.south) -- ([xshift=-5mm]maria.north);
\draw[flow] (mail) -- (ext);
\draw[flow] (ext.south) -- ([xshift=5mm]maria.north)
  node[lbl, pos=0.55, anchor=north] {schreibt direkt};
\draw[ref] (maria.west) -- (fs.east) node[lbl, midway] {Pfade};
\draw[flow] (ssm.south -| sqlite.north) -- (sqlite.north);
\draw[flow] (ssm.south -| csv.north) -- (csv.north);

% --- isoliertes Bewegungsmodell ---
\begin{scope}[on background layer]
  \node[draw=black!55, dash pattern=on 3pt off 2pt, line width=0.6pt,
        rounded corners=5pt, inner sep=7pt,
        fit=(hC)(ssm)(sqlite)(csv)] (isoC) {};
\end{scope}
\node[font=\scriptsize\itshape, text=black!60, anchor=north]
  at ($(isoC.south)+(0,-2pt)$) {keine Verbindung zu den\\[-2pt]übrigen Teilsystemen};

% --- Legende ---
\begin{scope}[shift={(3.0,-8.1)}, font=\scriptsize]
  \draw[flow] (0,0) -- (0.7,0);
  \node[anchor=west] at (0.8,0) {Datenzugriff};
  \draw[ref] (3.2,0) -- (3.9,0);
  \node[anchor=west] at (4.0,0) {Verweis über Dateipfad};
\end{scope}
\end{tikzpicture}
\caption[Ausgangszustand des Systems]{Ausgangszustand der drei Vorarbeiten.
  Kartenanwendung und Bildextraktion nutzen die gemeinsame MariaDB
  über getrennte Zugriffsklassen.}
\label{fig:ausgangszustand}
\end{figure}

\paragraph{Die vorgefundene Anwendung.} Die Kartenanwendung aus \cite{LorenzBA} hatte genau
zwei Seiten. Über ein Auswahlelement in der Seitenleiste ließ sich zwischen ihnen wechseln.
Jede Seite öffnete eine eigene Leaflet-Karte: entweder mit den Gerätestandorten in Japan
oder mit den Bewegungspfaden der Karpatenbären. Eine gemeinsame Darstellung beider
Datenbestände gab es nicht. Die japanische Ansicht bot bereits Funktionen zur
Verwaltung von Markern und Gebieten; die Bewegungssimulation war nicht eingebunden.

Das Bewegungsmodell aus \cite{BreuBA} lag
nur als Quelltext vor und war mit der Kartenanwendung nicht verbunden. Seine Ergebnisse
waren auf der Karte deshalb nicht zu sehen. Angewendet wurde es allein auf die
Karpatendaten, denn diese bestehen aus aufeinanderfolgenden Ortungen eines Tieres und lassen
sich in Schritte zerlegen. Die japanischen Daten geben dagegen Kamerastandorte und
Auslösezeitpunkte wieder und enthalten keine Bewegungsabfolge eines einzelnen Tieres
(Abschnitt~\ref{subsec:japandaten}).

Die Bildextraktion aus \cite{LelchatBA} war ein eigenständiger
Hintergrundprozess: Er holte neue Kameraaufnahmen aus einem E-Mail-Postfach und schrieb
sie in die Datenbank. Eine Kartendarstellung gehörte nicht zu seinem Umfang. Mit der
Kartenanwendung teilte er allein den Datenbestand.

\paragraph{Gemeinsame Datenhaltung.} Kartenanwendung und Bildextraktion nutzten
bereits dieselbe MariaDB (Abbildung~\ref{fig:ausgangszustand}), jedoch über getrennte
Zugriffsklassen: \texttt{WildlifeDatabase} und \texttt{DBManager}.
Das Bewegungsmodell nutzte eine eigene SQLite-Datenbank.

\paragraph{Daraus folgende Aufgabenstellung.} Aus diesem Zustand ergibt sich die zentrale
Aufgabe dieser Arbeit: Die beiden Datenbestände der Kartenanwendung und das bislang
getrennte Bewegungsmodell sollen in einem gemeinsamen System zusammengeführt werden, sodass
die Ergebnisse des Modells auf derselben Karte sichtbar werden wie die realen Daten.



\subsection{Zustand der Kartendarstellung}
\label{subsec:ausgangszustand-karte}

Beide Kartenansichten beruhten auf Leaflet, eingebunden über Folium. Die Art der Einbindung
richtete sich jeweils nach dem Zweck der Seite.

Die japanische Ansicht wurde über \texttt{st\_folium()} eingebunden und nutzte die
Folium-Erweiterung \texttt{Draw}. Damit ließen sich Marker, Rechtecke und Polygone auf der
Karte zeichnen. Die gezeichneten Objekte gelangten als GeoJSON zurück nach Python und wurden
dort auf Knopfdruck in die Datenbank geschrieben.

Die Karpaten-Ansicht diente der Darstellung der Bewegungsdaten und wurde über
\texttt{folium\_static()} eingebunden.\footnote{Zwei Aufrufe in
\texttt{carpathian\_bears.py}, jeweils mit fester Größe von 800$\times$600 Bildpunkten.}
Diese Funktion gibt die Karte als statisches HTML aus; sie besitzt, wie in
Abschnitt~\ref{subsec:leaflet} dargelegt, keinen Rückgabewert.

Jede der beiden Seiten erzeugte eine eigene Karteninstanz mit eigener Konfiguration. Beide
verwendeten den Standard-Renderer von Leaflet, also die Darstellung als Einzelelemente im
Dokumentbaum.

\begin{figure}[htbp]
\centering
\begin{subfigure}{0.48\linewidth}
  \includegraphics[width=\linewidth]{ausgangs-japan.png}
  \caption{Japan-Seite}
\end{subfigure}
\hfill
\begin{subfigure}{0.48\linewidth}
  \includegraphics[width=\linewidth]{ausgangs-karpaten.png}
  \caption{Karpaten-Seite}
\end{subfigure}
\caption[Kartenansichten im Ausgangszustand]{Die beiden Kartenansichten im
  Ausgangszustand. Nur die Japan-Seite nimmt Eingaben entgegen; die
  Werkzeugleiste am linken Kartenrand fehlt in der Karpaten-Ansicht.}
\label{fig:ausgangszustand-karten}
\end{figure}


% =====================================================================
\section{Datenbestände und Datenhaltung}
\label{sec:datenbestaende}

\subsection{GPS-Trajektorien der Karpaten-Studie}
\label{subsec:karpatendaten}

Der umfangreichste Datenbestand stammt aus der in Abschnitt~\ref{sec:motivation}
erwähnten Telemetriestudie an Europäischen Braunbären in den rumänischen Karpaten. Er liegt
als Datei \texttt{1\_bears\_RO.csv} vor und umfasst 59.997 Ortungen von 13~Tieren im
Zeitraum vom 9.~August 2011 bis zum 31.~Juli 2015.

Jede Zeile enthält neben Kennung und Zeitstempel die Koordinaten in EPSG:3844 (siehe
Abschnitt~\ref{subsec:crs}) sowie vier beschreibende Merkmale: Geschlecht, Altersklasse und
zwei Angaben zur Jahresphase. Tabelle~\ref{tab:karpaten-merkmale} zeigt deren Verteilung.

\begin{table}[htbp]
\centering
\caption{Verteilung der Merkmale im Karpaten-Datensatz ($n = 59.997$)}
\label{tab:karpaten-merkmale}
\small
\begin{tabular}{l r l r}
\toprule
\textbf{Jahresphase} & \textbf{Ortungen} & \textbf{Merkmal} & \textbf{Ortungen} \\
\midrule
Winter sleep              & 15.441 & männlich  & 36.315 \\
Den exit and reproduction & 13.199 & weiblich  & 23.682 \\
Forest fruits             & 15.091 & adult     & 41.052 \\
Hyperphagia               & 16.266 & juvenil   & 18.945 \\
\bottomrule
\end{tabular}
\end{table}

Zwei Eigenschaften dieses Bestands sind für die weitere Arbeit bedeutsam.

Erstens sind die Jahresphasen bereits Bestandteil der Rohdaten. Sie wurden nicht im Rahmen
dieser Arbeit oder einer Vorarbeit gebildet, sondern liegen als eigene Spalte vor. Die
Bezeichnungen folgen dem Jahreszyklus des Braunbären und sind damit fachlich und nicht
kalendarisch begründet.

Zweitens sind die Ortungen sehr ungleich auf die Tiere verteilt. Der am dichtesten
beobachtete Bär trägt 12.313 Ortungen bei, der am dünnsten beobachtete 1.024 - ein
Verhältnis von etwa zwölf zu eins. Auswertungen über alle Tiere hinweg werden dadurch von
wenigen Individuen dominiert. Dieser Umstand ist bei der Bewertung der in
Kapitel~\ref{chap:implementierung} vorgenommenen Messungen zu berücksichtigen.



\subsection{Kameradaten der Region Aizu}
\label{subsec:japandaten}

Der japanische Datenbestand ist strukturell anders und deutlich kleiner. Er besteht aus
20~Gerätestandorten im Aizu-Becken (Präfektur Fukushima) sowie den dort erfassten
Kameraaufnahmen. Die bekannten Standorte sind anhand ihrer
Gerätekennungen vier Nummernkreisen von
\texttt{001xxx} bis \texttt{004xxx} zugeordnet.



Nach der Bereinigung des Datenbestands stehen
49 reale Fotozeitpunkte von insgesamt
18 Gerätekennungen für die Auswertung zur Verfügung.
Die Standortdatenbank enthält 33 Gerätekennungen.
Davon verfügen 20 Geräte über gültige Koordinaten
und werden im Japan-Layer der Hauptkarte dargestellt.

Für die übrigen 13 Gerätekennungen liegen keine
gültigen Standortkoordinaten vor. Acht dieser
Geräte besitzen gespeicherte Kameraaufnahmen,
die über die Seitenleiste angezeigt und für
den zeitlichen Vergleich verwendet werden können.
Eine räumliche Darstellung dieser Geräte ist
ohne Standortkoordinaten nicht möglich.


Festzuhalten ist ferner, was über die Geräte \emph{nicht} gespeichert ist. Die Gerätetabelle
führt je Eintrag nur eine Kennung, den Namen des Bildverzeichnisses und die Koordinaten.
Die Bauart des Geräts oder die Art der Abschreckung ist nicht hinterlegt. Auswertungen, die
zwischen Gerätetypen unterscheiden, sind auf dieser Datengrundlage folglich nicht möglich.

\subsection{Datenhaltung im Ausgangszustand}
\label{subsec:datenhaltung}

Welches Teilsystem auf welche Ablage zugriff, zeigt Abschnitt~\ref{subsec:systemuebersicht}.
Hier geht es um die zentrale MariaDB-Instanz selbst.

Das Schema umfasste elf Tabellen: die Gerätestandorte, die von Nutzenden gezeichneten Marker
und Flächen, die Bildmetadaten, die Karpaten-Ortungen sowie einige abgeleitete Tabellen zu
Aktionsräumen und Tagesdistanzen. Für die Integration sind drei Eigenschaften relevant.

\begin{enumerate}
    \item In mehreren Tabellen sind Zeitstempel als
    Zeichenketten statt als Datumstyp gespeichert.
    Zeitliche Filter sind dadurch vom gespeicherten
    Format abhängig; Datumsfunktionen und die Prüfung
    ungültiger Werte sind erschwert.

    \item Die Koordinaten liegen in einfachen
    Gleitkommaspalten vor. Räumliche Datentypen kommen
    nicht zum Einsatz. Wer nur den sichtbaren
    Kartenausschnitt laden will, muss die Spalten
    einzeln vergleichen.

    \item Die Kameraaufnahmen liegen in zwei Tabellen.
    \texttt{image\_metadata} speichert je Aufnahme
    einen Dateipfad, das Bild selbst liegt im
    Dateisystem. \texttt{wildlife\_images} speichert
    die Bilddaten direkt in der Datenbank. Der Grund
    dafür ist dokumentiert: Für die Anbindung der
    Bildextraktion wurde die Ablage von Dateipfaden
    auf Bilddaten umgestellt \cite[S.~21]{LelchatBA},
    die ältere Tabelle blieb daneben bestehen. Eine
    gemeinsame Abfrage aller Aufnahmen muss daher
    beide Speicherformen berücksichtigen.
\end{enumerate}

Bei den Karpaten-Ortungen ist zudem auf mehrfach vorhandene Datensätze zu achten.
Der vorhandene Importcode kann mehrere CSV-Bestände zusammenführen, während die Tabelle
\texttt{bears\_tracking} keinen eindeutigen Schlüssel aus Tier und Zeitpunkt besitzt.
Mehrfach eingetragene Ortungen können die Darstellung und die Berechnung von
Bewegungskennzahlen beeinflussen.

% =====================================================================
\section{Übernommene Modelle und ihre Grenzen}
\label{sec:uebernommene-modelle}

\subsection{Bewegungssimulation}
\label{subsec:simulationsmodell}

\paragraph{Aufbau des Modells.}

Das Bewegungsmodell aus \cite{BreuBA} setzt das Prinzip der Step-Selection-Funktionen um
(Abschnitt~\ref{subsec:bewegungsoekologie}). Die Kandidatenerzeugung liegt in der Klasse
\texttt{StepSelectionModel}. Die Klasse \texttt{StepTrajectorySet} bildet Segmente aus den realen Trajektorien;
die Trennung nach Jahresphasen erfolgt über \texttt{SeasonalTrajectoryProcessor}.
\texttt{StepSelectionModel} nutzt diese vorbereiteten Daten.

Die Schrittlänge kommt aus einer Tabelle mit Geschwindigkeitswerten, die in der
mitgelieferten SQLite-Datei liegt. Die Tabelle hat 384~Einträge, gegliedert nach Geschlecht,
Altersklasse, Jahresphase und Stunde des Tages: zwei Geschlechter, zwei Altersklassen, vier
Jahresphasen und 24~Stunden. Je Eintrag stehen ein Mindest-, ein Höchst- und ein Medianwert
der Geschwindigkeit, wobei diese drei Werte bei 284 der 384~Einträge übereinstimmen.

Die Alternativschritte entstehen so: Für jeden realen Schritt sucht das Modell die mediane
Geschwindigkeit zu Geschlecht, Alter, Jahresphase und Stunde und rechnet daraus die
Schrittlänge. Dann bildet es $k$~Alternativen. Deren Richtung wird zufällig und
gleichverteilt über den ganzen Kreis gezogen. Je Alternative kommt ein Wahrnehmungspunkt in
200~Metern Entfernung dazu, an dem sich später Geländeeigenschaften auswerten lassen. Der
reale Schritt wird derselben Liste angefügt und als solcher markiert.

\paragraph{Ausgabe des Modells.}

Die Ausgabe ist eine Liste von Listen: je realem Schritt eine Menge von Kandidaten, jeder
mit Start- und Endpunkt, Wahrnehmungspunkt, Richtung, Länge und der Markierung, ob er der
reale Schritt ist.

Das Modell stellt damit mögliche Schritte bereit. Eine Bewertungsfunktion zur Auswahl
einzelner Kandidaten und zur Bildung zusammenhängender simulierter Pfade gehört nicht
zu diesem Teil der Vorarbeit. Für die Ziele dieser Arbeit wird diese Funktion ergänzt.

\paragraph{Erweiterungsbedarf für die Bewegungsanalyse.}
\label{subsubsec:limitationen-simulation}

Für die geplanten Untersuchungen sind vier Aspekte relevant.

\textbf{Richtungstreue.} Die Richtungen sind gleichverteilt über den ganzen Kreis. Die
Richtungspersistenz realer Tierbewegung (Abschnitt~\ref{subsec:bewegungsoekologie}) ist
darin nicht enthalten. Eine ungewichtete Auswahl der Kandidaten würde daher keinen
korrelierten Zufallspfad ergeben.

\textbf{Schrittlänge.} Alle Alternativen eines Schritts haben dieselbe Länge, abgeleitet aus
dem Median. Mindest- und Höchstwert gehen nicht ein, die Streuung der Schrittlängen damit
auch nicht - bei den oben genannten 284~Einträgen gibt es ohnehin keine.

\textbf{Zeitbezug der Schrittlänge.} Die Schrittlänge wird als Strecke einer Stunde
berechnet, unabhängig vom tatsächlichen Abstand zwischen zwei Ortungen. Bei längeren
Messintervallen ist deshalb eine zeitliche Anpassung der Schrittlänge sinnvoll.


\textbf{Abhängigkeiten.} Eine eigene Klasse fragt
Geländemerkmale über OpenStreetMap ab, ist
aber nicht an die Schrittauswahl angebunden.
Die ursprüngliche Implementierung führt zudem
eine HDBSCAN-basierte Clusteranalyse der
GPS-Ortungen durch. Dabei werden räumliche
Cluster und geometrische Bezugspunkte für
Trajektoriensegmente bestimmt.

Diese Ergebnisse werden bei der Erzeugung
alternativer Bewegungsschritte jedoch nicht
verwendet. Da die benötigten kompilierten
Bibliotheken auf dem Entwicklungsrechner
nicht zuverlässig geladen werden konnten,
wurde die Schritterzeugung ohne diese
nicht erforderliche Clusteranalyse umgesetzt.


\subsection{Bildextraktion}
\label{subsec:bildextraktion}

Die dritte Vorarbeit \cite{LelchatBA} übernimmt neue Kameraaufnahmen automatisch. Die
Kameras schicken ihre Aufnahmen als E-Mail-Anhang. Das System holt diese Anhänge ab, wertet
sie aus und trägt sie in die Datenbank ein.

\paragraph{Ablauf.}

Das System läuft als eigener Prozess, getrennt von der Kartenanwendung. Es fragt ein
Postfach über die Microsoft Graph API ab und meldet sich dabei per OAuth~2.0 als Anwendung
an, nicht als Benutzer. Die Abfrage läuft in einer Dauerschleife alle zehn Sekunden. Je
Durchlauf holt das System bis zu 50~ungelesene Nachrichten. Der Filter auf einen festen
Absender greift danach im Programm, nicht schon in der Abfrage.

Aus jedem Anhang im JPEG-Format gewinnt das System Information auf zwei Wegen:

\begin{itemize}
    \item \textbf{Aus dem Bild} liest die Texterkennung \texttt{easyOCR} den eingeblendeten
    Text. Daraus holt ein regulärer Ausdruck den Zeitstempel.
    \item \textbf{Aus dem Dateinamen} kommt die Gerätekennung. Das System kennt dafür zwei
    Namensformate. Trägt eine Aufnahme keinen eingeblendeten Zeitstempel - was vorkommt
    \cite[S.~25]{LelchatBA} --, stammt die Zeit ersatzweise aus dem Dateinamen. Das gelingt
    nur bei einem der beiden Formate, weil nur dieses ein Datum in auswertbarer
    Schreibweise enthält.
\end{itemize}

Danach schreibt das System die Aufnahme mit Gerätekennung und Zeitstempel in die Datenbank.
Ein Merkmal hält fest, ob beides bestimmt werden konnte. Ist die Gerätekennung in der
Gerätetabelle noch nicht vorhanden, legt das System dort einen neuen Eintrag an. Dieser
Eintrag erhält zunächst die Koordinaten 0/0.

\paragraph{Einordnung der Artbestimmung.}

Die Tierart wird nicht von diesem System bestimmt. An den Kameras im Gelände stehen
Raspberry-Pi-Rechner, die die Aufnahmen auswerten, Rahmen mit Klassenname und
Konfidenzwert in das Bild zeichnen und erst dann versenden \cite[S.~8]{LelchatBA}. Die
Kamera selbst blendet nur Gerätenummer, Zeitpunkt und Temperatur ein. Die Texterkennung
liest also die Beschriftung dieser Rahmen: Steht darin das Wort \enquote{bear}, wird die
Aufnahme übernommen, andernfalls verworfen \cite[S.~20]{LelchatBA}.

Die Artangabe stammt damit aus einem vorgeschalteten Modell, dessen Fehlerrate nicht bekannt
ist; Fehlerkennungen sind dokumentiert \cite[S.~24]{LelchatBA}. Sie wird, wie in
Abschnitt~\ref{sec:zielsetzung} festgelegt, als gegeben vorausgesetzt und in dieser Arbeit
weder geprüft noch verbessert.

\paragraph{Einbindung in die Kartenanwendung.}

Die Bildextraktion nutzt für den Postfachzugriff Berechtigungen, die vom IT-Zentrum
bereitgestellt werden müssen. Die Verarbeitung läuft unabhängig von der Kartenanwendung.
Beide Systeme teilen die Datenbank; eine Anzeige neu eingegangener Aufnahmen war in der
ursprünglichen Kartenoberfläche nicht vorgesehen.

% =====================================================================

\section{Anforderungen}
\label{sec:anforderungen}

Aus der Analyse der bestehenden Systeme und den Zielen der Arbeit
ergeben sich die folgenden Anforderungen. Im Mittelpunkt stehen die
gemeinsame Visualisierung der vorhandenen Daten sowie die Erweiterung
der Anwendung um Bewegungsanimationen und Vergleichsmöglichkeiten.

Für die Kartenvisualisierung ist MapLibre GL JS als technische
Grundlage vorgegeben. Die bisherigen Darstellungen mit
Leaflet und Folium werden dafür durch eine MapLibre-basierte Kartenansicht ersetzt.

\subsection{Funktionale Anforderungen}

\begin{description}

    \item[F1 - Gemeinsame Kartenansicht.]
    Die Daten der bestehenden Teilsysteme sollen auf einer gemeinsamen
    Karte dargestellt werden. Einzelne Datenebenen müssen unabhängig
    voneinander ein- und ausgeblendet werden können.

    \item[F2 - Datenfilterung.]
    Die GPS-Daten sollen nach Bär, Zeitraum und Jahresphase gefiltert
    werden können. Verschiedene Kartenansichten sollen zur Verfügung stehen.

    \item[F3 - Karteninteraktion.]
    Eigene Marker und Gebiete sollen direkt auf der Karte angelegt,
    gespeichert und wieder entfernt werden können.

    \item[F4 - Bewegungsanimation.]
    Reale und simulierte Bewegungen sollen zeitlich synchron
    dargestellt werden. Zeitraum und Abspielgeschwindigkeit sollen
    einstellbar sein.

    \item[F5 - Zusammenhängende Bewegungssimulation.]
    Aus den einzelnen möglichen Bewegungsschritten sollen
    zusammenhängende simulierte Pfade erzeugt werden können.

    \item[F6 - Simulation im japanischen Untersuchungsgebiet.]
    Das Bewegungsmodell aus den Karpaten soll auf die japanischen
    Kamerastandorte übertragen werden, um mögliche Bewegungsverläufe
    auch ohne vorhandene GPS-Trajektorien darzustellen.

    \item[F7 - Statistischer Vergleich.]
    Die simulierten Kameradetektionszeiten sollen mit den tatsächlichen
    Fotoauslösungen hinsichtlich ihrer zeitlichen Verteilung
    verglichen werden können.

    \item[F8 - Dateneinsicht und neue Aufnahmen.]
    Die verwendeten Kartendaten sollen einsehbar sein. Neu
    eingegangene Kameraaufnahmen sollen in der Anwendung
    kenntlich gemacht werden.

\end{description}

\subsection{Nicht-funktionale Anforderungen}

\begin{description}

    \item[N1 - Darstellungsleistung.]
    Die Karte soll auch bei größeren Datenbeständen und mehreren
    gleichzeitig sichtbaren Ebenen bedienbar bleiben.

    \item[N2 - Wartbarkeit.]
    Die Anwendung soll modular aufgebaut sein, damit bestehende
    Komponenten weiterverwendet und erweitert werden können.

    \item[N3 - Datenzugriff.]
    Die vorhandenen Datenbestände sollen über einheitliche
    Schnittstellen zugänglich sein. Für gemeinsam verwaltete
    Daten wird MariaDB verwendet.

    \item[N4 - Fehlerbehandlung.]
    Beim Ausfall einzelner Datenquellen sollen die übrigen Funktionen
    möglichst weiter nutzbar bleiben. Fehlende Verbindungen sollen
    erkennbar sein.

    \item[N5 - Reproduzierbarkeit.]
    Simulationen und statistische Auswertungen sollen bei gleichen
    Eingabedaten und Parametern reproduzierbar sein.

\end{description}


% ====================================================================
\chapter{Implementierung}
\label{chap:implementierung}

Dieses Kapitel beschreibt die technische Realisierung der integrierten Plattform.
Schwerpunkte bilden die Konzeption der interaktiven Kartenkomponente, die Konsolidierung
der heterogenen Teilsysteme und Datenbestände, die synchronisierte
Bewegungsanimation sowie die methodische Übertragung und Validierung des
Bewegungsmodells im japanischen Untersuchungsraum.

% =====================================================================


\section{Architektur und Integration der Anwendung}
\label{sec:architektur-integriert}

Die entwickelte Anwendung führt die vorhandenen Datenbestände und das Bewegungsmodell in einer gemeinsamen Streamlit-Oberfläche zusammen. Dabei bleiben die Datenaufbereitung, die Berechnung von Bewegungen und deren Darstellung voneinander getrennt. Abbildung~\ref{fig:architektur-integriert} zeigt den Aufbau des Systems.


\begin{figure}[htbp]
\centering

\begin{tikzpicture}[
    font=\footnotesize,
    >=Latex,
    node distance=5mm,
    view/.style={
        draw,
        rounded corners=2pt,
        align=center,
        text width=3.3cm,
        minimum height=10mm,
        fill=blue!5
    },
    process/.style={
        draw,
        rounded corners=2pt,
        align=center,
        text width=11.4cm,
        minimum height=10mm,
        fill=gray!8
    },
    data/.style={
        draw,
        align=center,
        text width=3.3cm,
        minimum height=10mm
    },
    flow/.style={->, thick},
    external/.style={->, dashed, thick}
]

% Drei Ansichten der Streamlit-Anwendung
\node[view] (main) at (0,0)
    {Hauptkarte};

\node[view] (carpathian) at (4.1,0)
    {Karpaten-Experiment};

\node[view] (japan) at (8.2,0)
    {Japan-Experiment};

% Umrahmung der Streamlit-Anwendung
\node[
    draw,
    rounded corners=4pt,
    inner sep=7mm,
    fit=(main)(carpathian)(japan),
    label={[font=\small\bfseries]above:Streamlit-Anwendung}
] (app) {};

% Gemeinsame Verarbeitung
\node[process] (processing) at (4.1,-2.0)
    {Datenaufbereitung, Bewegungsmodelle
     und statistische Auswertung};

% Datenquellen
\node[data] (mariadb) at (0,-4.0)
    {MariaDB\\Geräte, Bilder, GPS-Daten};

\node[data] (csv) at (4.1,-4.0)
    {CSV-Dateien\\GPS-Trajektorien};

\node[data] (sqlite) at (8.2,-4.0)
    {SQLite\\Geschwindigkeiten, Geländedaten};

% Externer Prozess
\node[data] (extractor) at (0,-5.7)
    {Bildextraktion\\externer Prozess};

% Aufrufbeziehungen
% Direkte Pfeile zur Verarbeitung
\draw[flow] (main.south) -- (main.south |- processing.north);
\draw[flow] (carpathian.south) -- (processing.north);
\draw[flow] (japan.south) -- (japan.south |- processing.north);

% Zugriff auf Datenquellen
% Direkte Pfeile zu den Datenquellen
\draw[<->, thick] (mariadb.north) -- (mariadb.north |- processing.south);
\draw[<->, thick] (csv.north) -- (processing.south);
\draw[<->, thick] (sqlite.north) -- (sqlite.north |- processing.south);

% Bildextraktion schreibt in MariaDB
\draw[external] (extractor.north)
    -- (mariadb.south);

\end{tikzpicture}

\caption{Vereinfachte Architektur der integrierten
Anwendung. Die drei Ansichten nutzen gemeinsame
Daten- und Verarbeitungsfunktionen. Die
gestrichelte Verbindung kennzeichnet die
derzeit nicht aktive Bildextraktion.}

\label{fig:architektur-integriert}

\end{figure}


Die Benutzeroberfläche umfasst drei Ansichten. Die \emph{Hauptkarte} stellt Kamerastandorte, Abschreckungsgeräte, GPS-Daten und Ergebnisse des übernommenen Bewegungsmodells als einzeln steuerbare Ebenen dar. Die \emph{Karpaten-Experimentseite} ermöglicht den zeitlichen Vergleich realer Bewegungen mit alternativen Schritten und zusammenhängenden simulierten Pfaden. Auf der \emph{Japan-Experimentseite} werden Bewegungswege ohne vorgegebene GPS-Trajektorie simuliert und virtuelle Kameraereignisse mit den tatsächlichen Auslösezeiten verglichen.

Die Ansichten greifen auf gemeinsame Daten und Berechnungsverfahren zurück, verwenden für die Darstellung jedoch eigene Kartenkomponenten. In der Hauptansicht bereitet Python die ausgewählten Daten als GeoJSON auf und übergibt sie an MapLibre. Die beiden Experimentseiten nutzen zusätzliche Simulations- und Animationsfunktionen. Die konkrete Kartenintegration wird in Abschnitt~\ref{sec:kartenvisualisierung} und die Bewegungsmodellierung in den folgenden Abschnitten behandelt.

Für die Datenhaltung werden bestehende Quellen weiterverwendet: MariaDB enthält unter anderem Geräte-, Bild- und GPS-Daten; lokale CSV- und SQLite-Dateien liefern Daten für die Bewegungsanalyse. Bei den Karpaten-Ortungen kann die Anwendung auf CSV-Dateien zurückgreifen, wenn die Datenbank nicht erreichbar ist. Die Bildextraktion bleibt ein eigenständiger Prozess außerhalb von Streamlit. Die Anwendung liest die gespeicherten Bildinformationen aus MariaDB und zeigt neue Einträge in der Seitenleiste an. Der automatische Postfachzugriff hängt von den dafür erforderlichen Berechtigungen ab und ist derzeit nicht nutzbar.



\section{Umsetzung der MapLibre-Karte}
\label{sec:kartenvisualisierung}

Für die gemeinsame Kartenvisualisierung wurde MapLibre GL JS
entsprechend der technischen Vorgabe in die bestehende
Streamlit-Anwendung eingebunden. Die Umsetzung trennt die
Datenaufbereitung in Python von der interaktiven Darstellung
im Browser.

\subsection{Einbindung und Datenübergabe}
\label{subsec:aufbau-karte}

Die Kartenkomponente wird über
\texttt{components.html()} in Streamlit eingebettet. Dazu
erzeugt die Funktion \texttt{render\_maplibre()} ein
HTML-Dokument, das MapLibre GL JS und die benötigten
Kartendaten enthält. Auf ein separates Frontend-Projekt
wird dadurch verzichtet.

Die Daten werden zuvor in Python aufbereitet und als
GeoJSON an die Kartenkomponente übergeben. Dazu gehören
die japanischen Geräte und Gebiete, die GPS-Daten der
Karpatenbären sowie die Ergebnisse der Bewegungssimulation.
Die Karte selbst stellt keine Datenbankabfragen, sondern
verarbeitet die bereits übergebenen Datenbestände.


Für die Darstellung der japanischen Kamerastandorte
werden nur Geräte mit gültigen Koordinaten als
Kartenpunkte übernommen. Geräte ohne bekannte
Standortinformationen erscheinen deshalb nicht
im Japan-Layer. Die zugehörigen Aufnahmen bleiben
jedoch in der Datenbank gespeichert und sind
über die Geräteauswahl in der Seitenleiste
weiterhin zugänglich.


\subsection{Ebenensteuerung und Kartendarstellung}

Die verschiedenen Datenbestände werden als eigene
GeoJSON-Quellen eingebunden und durch MapLibre-Ebenen
dargestellt. Punkte, Bewegungspfade und Gebiete erhalten
jeweils passende Darstellungsformen. Dadurch können die
einzelnen Datengruppen unabhängig voneinander ein- und
ausgeblendet werden.

Neben der Ebenensteuerung stehen unterschiedliche
Grundkarten zur Auswahl. Ein optionales digitales
Höhenmodell ermöglicht zudem eine dreidimensionale
Geländedarstellung. Beim Wechsel der Grundkarte werden
die eigenen Datenebenen nach dem Laden des neuen
Kartenstils erneut hinzugefügt, damit die ausgewählten
Informationen sichtbar bleiben.

Abbildung~\ref{fig:hauptkarte} zeigt die gemeinsame
Darstellung der Karpaten-GPS-Daten und der
Simulationsergebnisse.

\begin{figure}[htbp]
    \centering
    \includegraphics[width=0.75\linewidth]{karte-karpaten.png}
    \caption{Gemeinsame MapLibre-Karte mit GPS-Daten,
    Simulationsergebnissen und unabhängig steuerbaren
    Datenebenen.}
    \label{fig:hauptkarte}
\end{figure}

\subsection{Karteninteraktion und Zustandsverwaltung}
\label{sec:interaktion}

Die Steuerung der Ebenen, der Kartenansicht und der
Grundkarte erfolgt direkt im Browser. Da die Einbindung
über \texttt{components.html()} keinen regulären Rückkanal
zu Python bietet, bleiben diese Anzeigeänderungen
zunächst innerhalb der Kartenkomponente.

Damit die gewählte Ansicht bei einer erneuten Ausführung
von Streamlit erhalten bleibt, werden Mittelpunkt,
Zoomstufe, Kartenstil und sichtbare Ebenen in
URL-Parametern gespeichert. Beim erneuten Laden liest
Python diese Werte aus und übergibt sie wieder an
die Kartenkomponente.

Für das Anlegen eigener Marker und Gebiete ist dagegen
eine Rückgabe an Python erforderlich. Nach dem Zeichnen
werden die Koordinaten über URL-Parameter übertragen und
ein erneuter Seitenaufruf ausgelöst. Python liest die
übergebenen Werte aus und speichert das Objekt in MariaDB.
Anschließend werden die Parameter entfernt, um eine
erneute Verarbeitung zu vermeiden.

Das Löschen vorhandener Marker und Gebiete sowie die
Filterung der GPS-Daten erfolgen über Streamlit-Elemente
außerhalb der Karte. Damit werden die Karteninteraktion
und die Datenverwaltung innerhalb der gemeinsamen
Anwendung verbunden.






\section{Simulation und Animation der Karpatenbewegungen}
\label{sec:animation}

Für den Vergleich realer und simulierter Bewegungen wird das
Bewegungsmodell von Bräu in die Anwendung eingebunden.
Die ursprüngliche Schrittauswahl erzeugt zu den beobachteten
GPS-Bewegungen mehrere alternative Schritte. Darauf aufbauend
werden zusammenhängende simulierte Pfade berechnet und
gemeinsam mit der realen Trajektorie animiert. Abbildung~\ref{fig:karpaten-animation} zeigt die gemeinsame Darstellung einer realen
Trajektorie und zweier simulierter Bewegungsverläufe.


\begin{figure}[htbp]
    \centering
    \includegraphics[width=0.75\linewidth]{karpaten-animation.png}
    \caption{Zeitlich synchronisierte Darstellung einer realen
    GPS-Trajektorie (blau), zweier simulierter Bewegungswege
    (orange) und alternativer Bewegungsschritte (grau)
    im Karpaten-Experiment.}
    \label{fig:karpaten-animation}
\end{figure}


\subsection{Erzeugung zusammenhängender Bewegungswege}
\label{subsec:neuimplementierung}

Die Berechnung beginnt mit der Auswahl eines Bären und
eines Tages. Aus den zeitlich sortierten GPS-Positionen
werden aufeinanderfolgende Bewegungsschritte gebildet.
Für jeden Schritt werden alternative Richtungen erzeugt.
Die Schrittlänge ergibt sich aus den Geschwindigkeitsdaten
des bestehenden Modells und wird an den tatsächlichen
Zeitabstand zwischen zwei GPS-Positionen angepasst.

Während die alternativen Schritte jeweils von der realen
Ausgangsposition ausgehen, wird für einen simulierten Pfad
pro Zeitschritt ein Kandidat ausgewählt. Dessen Richtung
und Länge werden auf die zuletzt simulierte Position
übertragen. So entsteht ein zusammenhängender
Bewegungsverlauf, der sich von der realen Route
unterscheiden kann.

Die Kandidatenauswahl berücksichtigt auch die vorherige
Bewegungsrichtung. Dazu wird der mittlere Kosinus der
Drehwinkel verwendet, der für verschiedene Schrittlängen
aus den realen Karpaten-GPS-Daten bestimmt wurde.
Die Werte werden zwischen den untersuchten Schrittlängen
interpoliert und in einen Gewichtungsparameter
überführt. Die Auswahlwahrscheinlichkeit eines
Kandidaten ist proportional zu

\begin{equation}
    p_i \propto
    \exp\left(\kappa \cos(\theta_i)\right),
\end{equation}

wobei $\theta_i$ den Winkel zur vorherigen
Bewegungsrichtung und $\kappa$ den schrittabhängigen
Gewichtungsparameter bezeichnet. Dadurch wird
das Richtungsverhalten in Abhängigkeit von der
jeweiligen Schrittlänge berücksichtigt.

Da die Auswahl zufallsabhängig bleibt, werden mehrere
simulierte Pfade erzeugt. Diese stellen mögliche
Bewegungsverläufe unter den Modellannahmen dar und
sind keine Vorhersagen der tatsächlichen Route.
Ein Zufallsseed ist ein Startwert für den
Pseudozufallszahlengenerator, der die Reihenfolge
der erzeugten Zufallszahlen bestimmt. Bei gleichem
Seed und unveränderten Eingabeparametern werden
somit dieselben simulierten Bewegungsverläufe erzeugt.


\subsection{Synchronisierte Bewegungsanimation}
\label{subsec:synchron}

Für die Animation erhalten die simulierten Positionen
dieselben Zeitstempel wie die zugehörigen realen
GPS-Positionen. Dadurch lassen sich die Bewegungen
auf einer gemeinsamen Zeitachse darstellen.

Die Kartenkomponente zeigt die reale Route, die
alternativen Schritte und die simulierten Pfade in
unterschiedlichen Farben. Zwischen zwei aufeinanderfolgenden
Positionen wird die dargestellte Position anhand der
vergangenen Zeit linear interpoliert. Die Wiedergabe
kann gestartet, angehalten und schrittweise gesteuert
werden. Außerdem lassen sich die Geschwindigkeit
und der betrachtete Zeitraum einstellen.

Durch die gleichzeitige Darstellung wird sichtbar,
wie sich die simulierten Bewegungen von der
beobachteten Trajektorie unterscheiden. Die Animation
dient damit der Untersuchung des Modellverhaltens.
Eine Bewertung der Aussagekraft der Simulation
erfolgt in Kapitel~\ref{chap:diskussion}.


\section{Übertragung des Bewegungsmodells auf Japan}
\label{sec:japan-uebertragung}

Das Bewegungsmodell wird so eingesetzt,
dass es von einem Kamerastandort aus selbstständig aufeinanderfolgende
Positionen erzeugen kann. Die Schrittlängen stammen weiterhin aus den
Geschwindigkeitsdaten der Karpatenbären. Dabei handelt es sich um eine
Modellannahme, da die Daten von einer anderen Bärenart und aus einem
anderen Untersuchungsgebiet stammen.

\subsection{Simulation ohne vorgegebene Trajektorie}

Für jeden simulierten Zeitschritt wird die Schrittlänge anhand von
Geschlecht, Altersklasse, Jahreszeit und Tagesstunde bestimmt.
Für die Richtungswahl wird der Vollkreis in acht gleich große
Sektoren von jeweils $45^\circ$ unterteilt. Aus jedem Sektor
wird eine zufällige Richtung mit einer Abweichung von
höchstens $\pm 22{,}5^\circ$ zur Sektormitte erzeugt.
Dadurch schließen die Winkelbereiche lückenlos aneinander an,
sodass grundsätzlich jede Richtung als Kandidat vorkommen kann.

Die Anzahl von acht Sektoren ist ein festgelegter Modellparameter.
Sie ermöglicht eine Verteilung der Kandidaten über den gesamten
Kreis und begrenzt gleichzeitig den Berechnungsaufwand.
Bei aktivierter Geländegewichtung kommen weitere Richtungen
zu geeigneten Geländebereichen hinzu.

Die Richtungsauswahl erfolgt als gewichtete Zufallsentscheidung
und nicht durch die Auswahl des höchsten Bewertungswerts.
Dabei wird auch die vorherige Bewegungsrichtung berücksichtigt:
Richtungen, die den bisherigen Verlauf fortsetzen, erhalten
bei positivem Gewicht einen höheren Wert. Anders als bei der
Karpaten-Animation wird hierfür ein fester Gewichtungsparameter
verwendet. Die ausgewählte Richtung und Schrittlänge bestimmen
die nächste Position. Dieser Vorgang wird für die gewählte
Simulationsdauer wiederholt.

\subsection{Geländegewichtung und Bildung von Zielbereichen}


Da die Simulation ohne vorgegebene Trajektorie
keine Informationen über die tatsächlich
gewählten Bewegungswege besitzt, soll die
Geländegewichtung zusätzlich räumliche
Eigenschaften des Untersuchungsgebiets
bei der Richtungswahl berücksichtigen.

Optional werden Merkmale der Umgebung aus OpenStreetMap in die
Richtungswahl einbezogen. Beispielsweise erhalten Waldflächen
positive und stärker bebaute Flächen oder größere Straßen
negative Bewertungen. Diese Werte wurden als Modellparameter
festgelegt und nicht aus beobachteten japanischen
Bärenbewegungen geschätzt. Sie beschreiben daher angenommene
Präferenzen, keine nachgewiesene Lebensraumauswahl.

Die Suche nach geeigneten Geländebereichen erfolgt innerhalb
einer Umgebung von ungefähr 350 m um die aktuelle Position.
Damit können mehrere benachbarte Rasterzellen berücksichtigt
werden. Die Suchweite ist ein festgelegter Modellparameter
und stellt keine empirisch bestimmte Wahrnehmungsdistanz
der Bären dar.

Für die Bestimmung möglicher Zielrichtungen werden zunächst
Geländezellen mit einer ausreichend hohen positiven
Lebensraumbewertung ausgewählt. Benachbarte geeignete Zellen
werden über eine Acht-Nachbarschaft zu zusammenhängenden
Bereichen zusammengefasst. Anschließend wird für jeden
Bereich ein gewichteter Schwerpunkt berechnet:

\begin{equation}
    P = \frac{\sum_{i=1}^{n} w_i p_i}
             {\sum_{i=1}^{n} w_i}.
    \label{eq:japan-schwerpunkt}
\end{equation}

Dabei ist $p_i$ die Position einer Geländezelle und $w_i$
ihr positives Gewicht. Für die Gewichtung wird die
Bewertung oberhalb eines Mindestwerts verwendet und
um einen kleinen positiven Wert ergänzt. So können
auch Zellen am Mindestwert berücksichtigt werden.

Die Richtung vom aktuellen Standort zu jedem Schwerpunkt
wird als zusätzlicher Kandidat angeboten. Größere und
stärker gewichtete Bereiche erhalten dabei einen
höheren, nach oben begrenzten Zusatzwert.

Die Schwerpunktformel wird aus dem Verfahren von Geyer et al.
\cite{GeyerMelzerDuennweber2026} übernommen. Dort werden
nach der Trennung einzelner Fußabdrücke gewichtete
Positionen aus Sensormessungen berechnet. Hier wird
die Idee auf räumlich zusammenhängende Geländebereiche
übertragen. Das neuronale Segmentierungsverfahren
der Veröffentlichung wird dabei nicht verwendet.
Stattdessen werden die Bereiche anhand benachbarter
Rasterzellen gebildet.

Abbildung~\ref{fig:gelaendecluster} zeigt die einzelnen
Schritte an einem vereinfachten Geländeraster.

\begin{figure}[htbp]
    \centering
    \includegraphics[width=\linewidth]{04_ablauf_vier_schritte.png}
    \caption{Schematische Bildung geeigneter Zielbereiche:
    Auswahl und Bewertung von Geländezellen,
    Zusammenfassung benachbarter Zellen und
    Berechnung gewichteter Schwerpunkte.}
    \label{fig:gelaendecluster}
\end{figure}

\subsection{Bewertung und Auswahl der Bewegungsrichtung}

Die acht freien Richtungen und die Richtungen zu geeigneten
Geländebereichen bilden zusammen die Kandidatenmenge.
Bei eingeschalteter Geländegewichtung werden je Richtung
drei Geländepositionen entlang einer vorausliegenden
Strecke bewertet. Damit wird nicht nur das Schrittziel,
sondern auch die Umgebung entlang der möglichen
Bewegungsrichtung berücksichtigt.

Hinzu kommen gegebenenfalls der Zusatzwert für einen
geeigneten Bereich und die Bewertung der bisherigen
Bewegungsrichtung. Aus der Gesamtbewertung $S_i$
eines Kandidaten wird ein Auswahlgewicht berechnet:

\begin{equation}
    q_i = \exp(\beta S_i).
    \label{eq:japan-richtungsgewicht}
\end{equation}

Dabei bestimmt $\beta$, wie stark sich Unterschiede
zwischen den Bewertungen auf die Auswahl auswirken.
Die Gewichte werden für die zufällige Auswahl
der Kandidaten verwendet. Ein Kandidat mit höherer
Bewertung wird dadurch wahrscheinlicher, aber
nicht zwingend gewählt.

Für die wiederholte Geländeabfrage werden vorbereitete
OpenStreetMap-Daten in einem lokalen SQLite-Cache
gespeichert. Dadurch können die benötigten Rasterzellen
ohne erneute Netzwerkabfragen bereitgestellt werden.
Sind keine Geländeinformationen verfügbar, kann die
Simulation mit den freien Richtungen und der
Richtungstreue weiterlaufen.

Die Grenzen dieses vereinfachten Geländemodells,
insbesondere die festgelegten Geländegewichte und
die räumliche Suchweite, werden in
Kapitel~\ref{chap:diskussion} behandelt.



\section{Statistischer Vergleich der Kameradetektionszeiten}
\label{sec:tagesrhythmus}

Ein direkter Vergleich ist zwischen
simulierten und tatsächlichen Trajektorien nicht möglich.
Stattdessen wird untersucht, ob die simulierten Bären zu
ähnlichen Tageszeiten in die Nähe von Kameras gelangen,
zu denen auch reale Aufnahmen vorliegen.

Dafür werden aus den vorhandenen Kameradaten und den
simulierten Bewegungen zwei zeitliche Verteilungen gebildet.
Anschließend wird berechnet, wie stark sich diese
Verteilungen über einen Zeitraum von 24 Stunden überschneiden.
Abbildung~\ref{fig:aktivitaetsvergleich-ablauf} zeigt
den Ablauf des Vergleichs.

\begin{figure}[H]
    \centering
    \begin{tikzpicture}[
        >=Stealth,
        font=\small,
        node distance=0.9cm and 1.6cm,
        box/.style={
            draw,
            rounded corners=2pt,
            align=center,
            minimum height=1.0cm,
            text width=4.6cm,
            inner sep=4pt
        },
        widebox/.style={
            draw,
            rounded corners=2pt,
            align=center,
            minimum height=1.0cm,
            text width=6.3cm,
            inner sep=4pt
        },
        arrow/.style={->, thick}
    ]

    % linke Spalte
    \node[box] (real1) {Reale Kamerafotos\\Aufnahmezeitpunkte aus der Datenbank};
    \node[box, below=of real1] (real2) {Uhrzeiten der realen Aufnahmen};
    \node[box, below=of real2] (real3) {Zeitliche Dichtekurve\\der realen Detektionen};

    % rechte Spalte
    \node[box, right=3.0cm of real1] (sim1) {Simulierte Bewegungen\\300 Wege über jeweils 48 Stunden};
    \node[box, below=of sim1] (sim2)
    {Virtuelle Kameraeintritte\\beim Eintritt in einen 100\,m-Bereich};
    \node[box, below=of sim2] (sim3) {Zeitliche Dichtekurve\\der virtuellen Detektionen};

    % Mitte unten
    \node[widebox, below=1.2cm of $(real3)!0.5!(sim3)$] (delta)
    {Überlappungskoeffizient $\Delta$\\Gemeinsame Fläche beider Dichtekurven};

    \node[widebox, below=of delta] (boot)
    {1\,000 Bootstrap-Wiederholungen\\95\,\%-Konfidenzintervall};

    % Pfeile links
    \draw[arrow] (real1) -- (real2);
    \draw[arrow] (real2) -- (real3);

    % Pfeile rechts
    \draw[arrow] (sim1) -- (sim2);
    \draw[arrow] (sim2) -- (sim3);

    % Pfeile zur Mitte
    \draw[arrow] (real3.south) |- (delta.west);
    \draw[arrow] (sim3.south) |- (delta.east);

    % Pfeil nach unten
    \draw[arrow] (delta) -- (boot);

    \end{tikzpicture}
    \caption{Ablauf des statistischen Vergleichs zwischen realen und virtuellen Kameradetektionszeiten. Eigene Darstellung.}
    \label{fig:aktivitaetsvergleich-ablauf}
\end{figure}


\subsection{Erzeugung virtueller Kameradetektionen}
\label{subsec:monte-carlo}

Für den Vergleich werden mehrere Bewegungsverläufe mit
dem in Abschnitt~\ref{sec:japan-uebertragung}
beschriebenen Modell erzeugt. Dabei kommt ein
Monte-Carlo-Verfahren zum Einsatz. Das bedeutet,
dass die Simulation mehrfach mit zufällig gewählten
Ausgangsbedingungen durchgeführt wird.

In der Standardkonfiguration werden 300 Bewegungswege
über jeweils 48 Stunden berechnet. Zu Beginn jedes
Durchlaufs werden ein japanischer Kamerastandort,
Geschlecht, Altersklasse, Jahreszeit und Startstunde
zufällig ausgewählt. Die Startstunde wird gleichverteilt
aus den 24 Tagesstunden gezogen. Für diesen Vergleich
wird die Geländegewichtung deaktiviert.

Anschließend erzeugt das Modell für jede Stunde eine
neue Position. Für jede Position wird geprüft, ob sich
der simulierte Bär innerhalb eines Radius von 100\,m
um eine der bekannten Kameras befindet.

Als virtuelle Kameradetektion wird ein neuer Eintritt
in diesen Bereich gezählt. Befindet sich der Bär
über mehrere aufeinanderfolgende Zeitschritte im
Bereich derselben Kamera, wird nur der erste Eintritt
erfasst. Verlässt er den Bereich und betritt ihn
später erneut, entsteht eine weitere Detektion.
Dabei wird jede Kamera einzeln betrachtet, sodass
auch mehrere Eintritte zum selben Zeitpunkt möglich
sind. Kamerabereiche, in denen sich der Bär bereits
an seiner Startposition befindet, werden zu Beginn
nicht als Detektion gezählt.

Der Radius von 100\,m ist eine vereinfachte
Annahme für die Auswertung. Er entspricht nicht
der nachgewiesenen Erfassungsreichweite der
tatsächlich eingesetzten Kameras. Außerdem wird
nicht modelliert, ob eine Kamera bei einem
vorbeilaufenden Tier tatsächlich ausgelöst hätte.
Da die Positionen im Abstand von einer Stunde
berechnet werden, können kürzere Aufenthalte
im Kamerabereich unberücksichtigt bleiben.

Aus allen Simulationsdurchläufen entsteht eine
Sammlung virtueller Detektionszeiten. Diese wird
mit den Uhrzeiten der tatsächlichen Aufnahmen
verglichen. Da die Zahl der realen Aufnahmen je
Kamera gering ist, werden die verwendbaren
Aufnahmezeitpunkte aller Kameras gemeinsam
ausgewertet.



\subsection{Vergleich der zeitlichen Verteilungen}
\label{subsec:aktivitaetsvergleich}

Für den Vergleich reichen einzelne Uhrzeiten
allein nicht aus. Deshalb werden aus den realen
und virtuellen Detektionszeiten zunächst zwei
Dichtekurven berechnet. Dafür wird eine
Kerndichteschätzung verwendet.

Bei diesem Verfahren wird um jeden Aufnahmezeitpunkt
eine kleine Glockenkurve gelegt. Die einzelnen
Glockenkurven werden zusammengefasst, sodass
eine durchgehende Verteilung über den Tag entsteht.
Bereiche mit vielen nahe beieinanderliegenden
Detektionen erhalten dabei eine höhere Dichte.
Die verwendete Bandbreite beträgt 1,5 Stunden.

Da sich der Tagesrhythmus alle 24 Stunden wiederholt,
muss auch der Übergang über Mitternacht berücksichtigt
werden. Beispielsweise liegen 23:30 Uhr und
00:30 Uhr nur eine Stunde auseinander. Die
Implementierung verwendet deshalb eine periodisch
fortgesetzte Gauß-Kerndichteschätzung. Dabei
werden die Zeitpunkte zusätzlich um 24 Stunden
verschoben, damit am Tageswechsel keine
künstliche Unterbrechung entsteht.

Die beiden Dichtekurven werden jeweils so
normiert, dass ihre gesamte Fläche den Wert
eins besitzt. Somit können auch Stichproben
mit unterschiedlich vielen Detektionen
miteinander verglichen werden.

Zur Beschreibung der zeitlichen Übereinstimmung
wird der Überlappungskoeffizient $\Delta$
nach Ridout und Linkie
\cite{RidoutLinkie2009} verwendet:

\begin{equation}
    \Delta =
    \int_{0}^{24}
    \min\left(
        f_{\mathrm{real}}(t),
        f_{\mathrm{sim}}(t)
    \right)\,dt.
    \label{eq:aktivitaetsueberlappung}
\end{equation}

Dabei beschreibt $f_{\mathrm{real}}(t)$ die
geschätzte Dichte der realen Aufnahmezeiten
und $f_{\mathrm{sim}}(t)$ die geschätzte Dichte
der virtuellen Detektionszeiten zur Uhrzeit $t$.
Für jeden Zeitpunkt wird der kleinere Wert
der beiden Kurven berücksichtigt. Die
gemeinsame Fläche über den gesamten Tag
ergibt den Überlappungskoeffizienten.

Der Wert $\Delta$ liegt zwischen null und eins.
Ein Wert nahe null bedeutet, dass sich die beiden
zeitlichen Verteilungen nur wenig überschneiden.
Ein Wert nahe eins weist dagegen auf eine starke
Übereinstimmung hin. Die Fläche wird in der
Implementierung numerisch über kleine
Zeitabschnitte berechnet.

\paragraph{Bootstrap-Verfahren zur Abschätzung der Unsicherheit}

Da die Zahl der realen Kameradetektionen begrenzt ist,
kann der berechnete Überlappungskoeffizient von den
vorliegenden Aufnahmezeitpunkten beeinflusst werden.
Um diese statistische Unsicherheit abzuschätzen,
wird ein Bootstrap-Verfahren verwendet.

Dabei werden aus den realen und virtuellen
Detektionszeiten jeweils neue Stichproben mit
Zurücklegen gezogen. Für jede Stichprobe werden
die beiden Dichtekurven und der
Überlappungskoeffizient erneut berechnet.
Dieser Vorgang wird 1\,000-mal wiederholt.
Die Anzahl wurde als praktikabler Wert gewählt,
um die Streuung der Ergebnisse abzuschätzen,
ohne den Berechnungsaufwand zu stark zu erhöhen.

Zur Darstellung der Unsicherheit wird ein
95-\%-Bootstrap-Konfidenzintervall verwendet.
Dieses Konfidenzniveau ist in statistischen
Untersuchungen weit verbreitet und stellt
einen Kompromiss zwischen der Breite des
Intervalls und dem gewählten Konfidenzniveau dar.

Die 1\,000 berechneten Überlappungswerte werden
der Größe nach geordnet. Die untere Grenze
des Intervalls bildet das 2,5-\%-Perzentil,
die obere Grenze das 97,5-\%-Perzentil.
Dadurch liegen die mittleren 95\,\% der
Bootstrap-Ergebnisse zwischen diesen Grenzen.
Die jeweils äußersten 2,5\,\% werden bei der
Bestimmung des Intervalls nicht berücksichtigt,
weil ein zweiseitiges Intervall mit dem
Konfidenzniveau von 95\,\% verwendet wird.

Das Intervall beschreibt die statistische
Unsicherheit der geschätzten zeitlichen
Überlappung unter den Annahmen des
Bootstrap-Verfahrens. Es gibt dagegen
nicht die Genauigkeit des Bewegungsmodells an.

\paragraph{Einordnung des Vergleichs}

Der statistische Vergleich wird in seiner
Standardkonfiguration ohne zusätzliche
Geländegewichtung durchgeführt. Die Ergebnisse
zeigen daher zunächst, wie stark die mit dem
übertragenen Bewegungsmodell erzeugten
Detektionszeiten mit den realen Aufnahmezeiten
übereinstimmen.

Eine hohe zeitliche Überlappung bedeutet jedoch
nicht, dass die simulierten Bären dieselben
Wege wie die tatsächlich beobachteten Tiere
zurücklegen. Ebenso wird kein eigener
Tagesaktivitätsrhythmus japanischer Bären
aus Beobachtungsdaten abgeleitet. Die Analyse
dient deshalb als erster statistischer Vergleich
der zeitlichen Muster und nicht als vollständige
biologische Validierung des Bewegungsmodells.
Die Ergebnisse und ihre Grenzen werden in
Kapitel~\ref{chap:diskussion} betrachtet.



% ====================================================================
\chapter{Ergebnisse und Diskussion}
\label{chap:diskussion}

Dieses Kapitel stellt die Ergebnisse der
Funktionsprüfungen und der Simulationsuntersuchungen
vor. Dabei werden die technischen Funktionen
bewertet und die Aussagekraft sowie die Grenzen
der Simulationen eingeordnet.


\section{Ergebnisse der Anwendungsintegration}
\label{sec:ergebnisse-integration}

Die Funktionsprüfung zeigt, dass die zuvor getrennten
Datenbestände in einer gemeinsamen Kartenanwendung
dargestellt werden können. Auf der Hauptkarte lassen sich
die GPS-Daten aus den Karpaten, die Ergebnisse der
Bewegungssimulation und die japanischen Kamerastandorte
als eigenständige Datenebenen einblenden.

Bei den japanischen Kameras ist die
Kartendarstellung auf Geräte mit bekannten
Standortkoordinaten beschränkt. Aufnahmen
von Geräten ohne diese Angaben sind
weiterhin über die Seitenleiste zugänglich.

Bei der Erprobung konnten die GPS-Daten und die
Simulationsdaten gleichzeitig angezeigt werden. Die
einzelnen Ebenen ließen sich unabhängig voneinander
aktivieren und deaktivieren. Auch der Wechsel zwischen
verschiedenen Kartenstilen funktionierte, ohne dass
die ausgewählten Datenebenen verloren gingen.
Abbildung~\ref{fig:ergebnis-hauptkarte} zeigt die
gemeinsame Kartenansicht mit eingeblendeten
Karpaten-GPS- und Simulationsdaten.

\begin{figure}[htbp]
    \centering
    \includegraphics[width=\linewidth]{ergebnis-hauptkarte.png}
    \caption{Integrierte MapLibre-Karte mit gleichzeitig
    dargestellten GPS-Daten aus den Karpaten und
    Ergebnissen der Bewegungssimulation. Die Datenebenen
    können über die Kartensteuerung einzeln ausgewählt werden.}
    \label{fig:ergebnis-hauptkarte}
\end{figure}

Auch die Filterfunktionen wurden überprüft. Die Auswahl
eines Bären, eines Zeitraums oder einer Jahresphase
veränderte die angezeigten Daten entsprechend. Die
zugehörigen Datensätze konnten zusätzlich in den Tabellen
unterhalb der Karte eingesehen werden.

Das Ein- und Ausblenden der Kartenebenen erfolgt direkt
im Browser und erfordert keinen erneuten Durchlauf der
Streamlit-Anwendung. Bei der Änderung von Datenfiltern
wird die Anwendung dagegen weiterhin neu ausgeführt.
Durch Formulare und die Zwischenspeicherung bereits
geladener Daten werden unnötige Datenabfragen vermieden.
Während der praktischen Erprobung wurde die Bedienung
beim Filterwechsel als flüssiger wahrgenommen als in der
vorherigen Anwendung. Eine systematische Messung der
Ladezeiten wurde jedoch nicht durchgeführt.

Die Funktionen zum Anlegen und Speichern eigener Marker
und Flächen wurden ebenfalls erfolgreich getestet.
Die gezeichneten Koordinaten werden von der Kartenkomponente
an die Streamlit-Seite übermittelt und anschließend in
der Datenbank gespeichert. Nach dem Speichern wird die
gesamte Seite neu geladen. Die gespeicherten Einträge
können danach wieder auf der Karte angezeigt werden.

Zusätzlich bietet die Anwendung eine dreidimensionale
Geländedarstellung. Diese wurde erfolgreich aufgerufen
und ermöglicht eine räumlich anschauliche Betrachtung
des Geländes. Allerdings traten während der Erprobung
gelegentlich Verzögerungen beim Laden des Höhenmodells
sowie eine stockende Kartenbedienung auf. Die
zweidimensionale Darstellung blieb davon unabhängig
nutzbar.

Insgesamt bestätigt die Funktionsprüfung die
Zusammenführung der zuvor getrennten Kartenansichten.
Die zentralen Datenbestände können innerhalb einer
gemeinsamen Oberfläche betrachtet und bedient werden.
Einschränkungen zeigten sich hauptsächlich bei der
Darstellungsleistung der optionalen 3D-Ansicht.




\section{Ergebnisse der Karpaten-Simulation}
\label{sec:ergebnisse-karpaten}

\subsection{Synchronisierte Bewegungsdarstellung}

Die Karpaten-Simulation wurde anhand realer GPS-Daten
besenderter Braunbären überprüft. Dabei wurde untersucht,
ob sich die aufgezeichneten Bewegungen gemeinsam mit
den vom Modell erzeugten Bewegungsverläufen darstellen
lassen.

Für die Erprobung wurde unter anderem Bear12 am
05.07.2013 ausgewählt. Die Anwendung zeigte die reale
GPS-Bewegung und drei simulierte Bewegungsverläufe
innerhalb derselben Kartenansicht. Alle dargestellten
Bewegungen verwenden eine gemeinsame Zeitachse.

Beim Abspielen bewegten sich der reale Bär und die
simulierten Bären entlang ihrer jeweiligen Trajektorien.
Auch das Anhalten, das schrittweise Navigieren sowie
das manuelle Verschieben der Zeitleiste funktionierten
bei der praktischen Erprobung.

Die gemeinsame Zeitachse ermöglicht es, die Positionen
des realen Bären und der simulierten Bären zu einem
beliebigen dargestellten Zeitpunkt gegenüberzustellen.
Die Simulation bleibt dabei zeitlich an die realen
Messpunkte gebunden, ohne deren räumlichen Verlauf
übernehmen zu müssen.

\subsection{Vergleich der Bewegungsverläufe}

Die erzeugten Bewegungsverläufe zeigen, dass aus
den Kandidatenschritten des Bewegungsmodells
zusammenhängende simulierte Pfade gebildet werden
können. Die simulierten Bären beginnen am
Ausgangspunkt der realen Trajektorie. Im weiteren
Verlauf können sich ihre Bewegungsrichtungen
und Positionen voneinander unterscheiden.

Abbildung~\ref{fig:ergebnis-karpaten-animation}
zeigt ein Beispiel aus der Karpaten-Animation.
Die reale Bewegung ist blau dargestellt,
während die simulierten Bewegungsverläufe
in unterschiedlichen Farbtönen erscheinen.
Die zusätzlich eingeblendeten alternativen
Schrittpositionen zeigen mögliche Bewegungsrichtungen
für den jeweiligen GPS-Schritt.

\begin{figure}[htbp]
    \centering
    \includegraphics[
        width=0.95\linewidth,
        height=0.65\textheight,
        keepaspectratio
    ]{ergebnis-karpaten-animation.png}
    \caption{Gegenüberstellung der realen GPS-Bewegung
    von Bear12 und drei simulierten Bewegungsverläufen
    am 05.07.2013. Die Darstellung zeigt die Bewegungen
    zu einem gemeinsamen Zeitpunkt sowie die
    verfügbaren alternativen Schrittpositionen.
    Eigene Darstellung.}
    \label{fig:ergebnis-karpaten-animation}
\end{figure}

Die simulierten Pfade folgen nicht zwangsläufig
der tatsächlich aufgezeichneten Route. Dies
entspricht der vorgesehenen Funktionsweise:
Jeder simulierte Pfad stellt eine mögliche
Realisierung des Bewegungsmodells dar.
Er ist weder eine Rekonstruktion noch eine
Vorhersage der konkreten GPS-Trajektorie.

Die Darstellung mehrerer Realisierungen macht
sichtbar, dass aus denselben Ausgangsdaten
unterschiedliche Bewegungsverläufe entstehen
können. Die Entfernung zwischen realem und
simuliertem Bären lässt sich während der
Animation zusätzlich als mittlerer Abstand
ablesen. Dieser Wert beschreibt die räumliche
Abweichung zum jeweiligen Zeitpunkt, ist jedoch
kein unmittelbares Maß für die Güte des
Bewegungsmodells.



\section{Ergebnisse der Simulation in Japan}
\label{sec:ergebnisse-japan}

\subsection{Darstellung der simulierten Bewegungen}

Die Japan-Simulation wurde daraufhin überprüft,
ob das Bewegungsmodell auch ohne vorgegebene
GPS-Trajektorie zusammenhängende Bewegungsverläufe
erzeugen kann. Als Ausgangspunkt dient dabei
ein ausgewählter Kamerastandort.

Für die praktische Erprobung wurde unter anderem
die Kamera 002001 verwendet. Von diesem Standort
aus wurden sechs Bären gleichzeitig simuliert.
Die Anwendung stellte deren Bewegungsverläufe
mit unterschiedlichen Farben auf einer Karte dar.

Beim Abspielen konnten die Bewegungen über eine
gemeinsame Zeitleiste verfolgt werden. Die
simulierten Bären starteten am selben Standort,
bewegten sich anschließend jedoch auf
unterschiedlichen Wegen. Zusätzlich wurden die
Standorte der vorhandenen Kameras angezeigt.

Näherte sich ein simulierter Bär einer Kamera
innerhalb des eingestellten Erfassungsradius,
wurde dies in der Anwendung als simulierte
Kameraauslösung angezeigt. Für den ausgewählten
Startstandort konnten außerdem die Zeitpunkte
realer Fotoaufnahmen auf der Zeitleiste
eingeblendet werden.

Damit ließ sich die Simulation sowohl räumlich
als auch zeitlich betrachten. Da für Japan keine
durchgehenden GPS-Trajektorien einzelner Bären
vorliegen, konnten die erzeugten Wege jedoch
nicht direkt mit tatsächlich zurückgelegten
Routen verglichen werden.
Ein Beispiel der Kartenansicht ist in
Abbildung~\ref{fig:ergebnis-japan-gelaende}
dargestellt.


\subsection{Einfluss der Geländegewichtung}

Um den Einfluss der Geländegewichtung sichtbar
zu machen, wurden Simulationen mit deaktivierter
und aktivierter Geländegewichtung betrachtet.
Abbildung~\ref{fig:ergebnis-japan-gelaende}
zeigt beispielhafte Bewegungsverläufe aus
beiden Varianten.

\begin{figure}[p]
    \centering

    \begin{subfigure}{0.98\linewidth}
        \centering
        \includegraphics[
            width=\linewidth
        ]{ergebnis-japan-animation-no-gl.png}
        \caption{Ohne Geländegewichtung}
    \end{subfigure}

    \vspace{0.5cm}

    \begin{subfigure}{0.98\linewidth}
        \centering
        \includegraphics[
            width=\linewidth
        ]{ergebnis-japan-animation-with-gl.png}
        \caption{Mit Geländegewichtung}
    \end{subfigure}

    \caption{Beispielhafte Bewegungsverläufe
    von sechs simulierten Bären ausgehend von
    Kamera 002001. Dargestellt sind Simulationen
    ohne und mit Geländegewichtung.
    Die dreidimensionale Geländeansicht dient
    ausschließlich der Visualisierung.
    Eigene Darstellung.}
    \label{fig:ergebnis-japan-gelaende}
\end{figure}

Ohne Geländegewichtung werden die möglichen
Bewegungsrichtungen unabhängig von den
Landschaftsmerkmalen bewertet. Die Auswahl
berücksichtigt jedoch weiterhin die vorherige
Bewegungsrichtung.

Bei aktivierter Geländegewichtung werden
zusätzlich die Merkmale der Umgebung aus
OpenStreetMap einbezogen. Beispielsweise können
Richtungen zu geeigneten Waldflächen höhere
Auswahlgewichte erhalten. Dadurch verändern
sich die Wahrscheinlichkeiten der möglichen
Bewegungsrichtungen.

Die beiden Darstellungen veranschaulichen
die unterschiedlichen Einstellungen der
Simulation. Da die Bewegungswege zufallsabhängig
entstehen, reicht der Vergleich einzelner
Pfade jedoch nicht aus, um eine Verbesserung
des Bewegungsmodells nachzuweisen. Hierfür
müssen mehrere Simulationsläufe anhand
messbarer Merkmale ausgewertet werden.


\subsection{Vergleich der clusterbasierten Zielrichtungen}

Zusätzlich wurde untersucht, welchen Einfluss
die in Kapitel~\ref{sec:japan-uebertragung}
beschriebenen Zielrichtungen zu geeigneten
Geländebereichen auf die Bewegungen haben.

Ziel dieses Tests war es zu prüfen, ob die
zusätzlichen Richtungen die simulierten Bären
häufiger in Bereiche mit geeigneten
Geländemerkmalen führen.

Dafür wurden zwei Varianten verglichen.
Die erste Variante verwendete nur die freien
Richtungskandidaten mit Geländegewichtung.
In der zweiten Variante kamen zusätzlich
die Richtungen zu den gewichteten
Schwerpunkten geeigneter Geländebereiche
hinzu.

Für beide Varianten wurden jeweils
60 Simulationsläufe über 24 Stunden
durchgeführt. Dabei wurden dieselben
Kamerastandorte, Ausgangsbedingungen
und Zufallsseeds verwendet. So konnte
der Einfluss der zusätzlichen
Zielrichtungen untersucht werden.

Die Ergebnisse wurden anhand des Anteils
von Positionen mit bestimmten
OpenStreetMap-Merkmalen, der mittleren
Geländebewertung und der Nettoverlagerung
ausgewertet. Die Nettoverlagerung bezeichnet
die direkte Entfernung zwischen Start-
und Endposition eines simulierten Bären.

Tabelle~\ref{tab:ergebnis-japan-cluster}
zeigt die Mittelwerte der Testläufe.

\begin{table}[htbp]
    \centering
    \caption{Ergebnisse der Japan-Simulation
    ohne und mit zusätzlichen clusterbasierten
    Zielrichtungen. Beide Varianten verwenden
    die Geländegewichtung. Mittelwerte aus
    jeweils 60 Simulationsläufen über 24 Stunden.
    Eigene Auswertung.}
    \label{tab:ergebnis-japan-cluster}
    \small

    \begin{tabular}{L{5.0cm}rr}
        \toprule
        \textbf{Kennzahl}
        & \textbf{Ohne Cluster}
        & \textbf{Mit Cluster} \\
        \midrule

        Positionen mit Waldmerkmal
        & 31,9\,\%
        & 39,1\,\% \\

        Positionen mit Wohnstraßenmerkmal
        & 38,1\,\%
        & 42,7\,\% \\

        Mittlere Geländebewertung
        & $-0{,}37$
        & $-0{,}27$ \\

        Nettoverlagerung
        & 976\,m
        & 909\,m \\

        \bottomrule
    \end{tabular}
\end{table}

Mit den zusätzlichen Zielrichtungen stieg
der Anteil der Positionen mit Waldmerkmal
von 31,9\,\% auf 39,1\,\%. Das entspricht
einem Anstieg um 7,2 Prozentpunkte.
Auch die mittlere Geländebewertung
verbesserte sich von $-0{,}37$ auf
$-0{,}27$.

Gleichzeitig erhöhte sich der Anteil der
Positionen mit Wohnstraßenmerkmal von
38,1\,\% auf 42,7\,\%. Dieses Merkmal
beschreibt Wohnstraßen und ist nicht
mit bebauten Wohngebieten gleichzusetzen.
Die Ergebnisse zeigen somit, dass eine
stärkere Bevorzugung von Waldmerkmalen
nicht automatisch zu einer geringeren
Nähe zu allen Straßen führt.

Die mittlere Nettoverlagerung verringerte
sich von 976\,m auf 909\,m. Die
simulierten Bären entfernten sich
damit in der Variante mit zusätzlichen
Zielrichtungen im Durchschnitt etwas
weniger weit vom Ausgangspunkt.

Insgesamt zeigen die Testergebnisse,
dass die zusätzlichen Zielrichtungen
die räumliche Verteilung der
simulierten Positionen beeinflussen.
Unter den gewählten Modellannahmen
wurden häufiger Positionen mit
Waldmerkmal erreicht.

Daraus lässt sich jedoch nicht
ableiten, dass die erzeugten Wege
dem tatsächlichen Bewegungsverhalten
japanischer Bären entsprechen.
Die Geländegewichte wurden festgelegt
und nicht anhand japanischer GPS-Daten
bestimmt. Außerdem hängt die Bewertung
von den verfügbaren OpenStreetMap-Daten ab.

Eine weitere Möglichkeit zur Überprüfung
bietet der Vergleich der simulierten
Kameraauslösungen mit den tatsächlich
aufgezeichneten Fotozeitpunkten.
Diese statistische Auswertung wird
im folgenden Abschnitt betrachtet.




\section{Statistischer Vergleich der Kameradetektionszeiten}
\label{sec:ergebnisse-kamerazeiten}

\subsection{Durchführung der Auswertung}

Für den zeitlichen Vergleich wurden die tatsächlichen
Fotozeitpunkte der japanischen Kameras den virtuellen
Kameraeintritten der Simulation gegenübergestellt.
Dabei wurden die Kameras gemeinsam betrachtet,
da für einzelne Kamerastandorte nur wenige reale
Aufnahmen vorliegen.

Insgesamt standen 49 reale Fotozeitpunkte zur Verfügung.
Für die virtuellen Ereignisse wurden 300 Simulationsläufe
mit jeweils 48 Stunden Dauer durchgeführt. Jeder Lauf
begann an einem zufällig ausgewählten Kamerastandort
und zu einer zufälligen Tagesstunde. Auch Geschlecht,
Altersklasse und Jahreszeit wurden zufällig gewählt,
da diese Angaben für die fotografierten Tiere
nicht bekannt sind. Die Geländegewichtung war
für diesen Vergleich deaktiviert.

Als virtueller Kameraeintritt wurde gezählt,
wenn ein simulierter Bär bei den stündlich
berechneten Positionen neu in den Umkreis von
100\,m einer Kamera gelangte. Ein längerer Aufenthalt
im selben Kamerabereich erzeugte keine zusätzlichen
Eintritte. Der Start an einem Kamerastandort
wurde ebenfalls nicht als Eintritt gezählt.
Der Radius beschreibt dabei eine angenommene
räumliche Nähe und nicht die tatsächliche
Erfassungsreichweite der Kamera.

Die Anzahl der Simulationsläufe wurde gewählt,
um nicht nur einzelne zufällige Bewegungswege
zu betrachten. Die Dauer von 48 Stunden umfasst
zwei vollständige Tageszyklen. Für die
Reproduzierbarkeit wurde der Zufallsseed 42
verwendet. Mit diesen Einstellungen entstanden
247 virtuelle Kameraeintritte.

\subsection{Vergleich der Tageszeitverteilungen}

Für die realen Fotozeitpunkte und die virtuellen
Kameraeintritte wurden die Tageszeitverteilungen
mit dem in Abschnitt~\ref{sec:tagesrhythmus}
beschriebenen Verfahren geschätzt. Die verwendete
KDE-Bandbreite betrug 1,5 Stunden. Dabei handelt
es sich um einen festgelegten Glättungsparameter.
Die Berechnung berücksichtigt den Übergang
zwischen 24:00 und 00:00 Uhr.

Abbildung~\ref{fig:ergebnis-kamerazeiten-kde}
zeigt die beiden Verteilungen. Die schattierte
Fläche stellt den gemeinsamen Bereich der
Dichtekurven dar.

\begin{figure}[htbp]
    \centering
    \includegraphics[width=0.95\linewidth]{ergebnis-kamerazeiten-kde.pdf}
    \caption{Tageszeitverteilungen der 49 realen Fotozeitpunkte
    und der 247 virtuellen Kameraeintritte. Die gemeinsame Fläche
    unter den Dichtekurven entspricht dem Überlappungskoeffizienten.
    Eigene Auswertung.}
    \label{fig:ergebnis-kamerazeiten-kde}
\end{figure}

Für die gewählten Einstellungen ergab sich ein
Überlappungskoeffizient von

\begin{equation}
    \Delta = 0{,}522.
\end{equation}


Der Wert zeigt, dass sich die geschätzten
Tageszeitverteilungen teilweise überschneiden.
Er gibt jedoch nicht an, welcher Anteil der
simulierten Kameraereignisse tatsächlichen
Fotoaufnahmen entspricht. Verglichen wurden
ausschließlich die Verteilungen der Uhrzeiten,
nicht einzelne Ereignisse oder tatsächliche
Bewegungswege.


\subsection{Unsicherheit und Einordnung}

Zur Abschätzung der Unsicherheit wurde das
Bootstrap-Verfahren mit 1000 Wiederholungen
angewendet. Dabei wurden aus den realen und den
virtuellen Ereigniszeiten jeweils neue Stichproben
mit Zurücklegen gezogen und der Überlappungskoeffizient
erneut berechnet. Abbildung~\ref{fig:ergebnis-kamerazeiten-bootstrap}
zeigt die Verteilung dieser Werte.

\begin{figure}[htbp]
    \centering
    \includegraphics[width=0.95\linewidth]{ergebnis-kamerazeiten-bootstrap.pdf}
    \caption{Verteilung des Überlappungskoeffizienten aus
    1000 Bootstrap-Wiederholungen. Die senkrechte durchgezogene
    Linie zeigt den Wert der ursprünglichen Auswertung.
    Die gestrichelten Linien markieren das
    95\,\%-Bootstrap-Intervall.
    Eigene Auswertung.}
    \label{fig:ergebnis-kamerazeiten-bootstrap}
\end{figure}

Das berechnete 95-\%-Bootstrap-Intervall
lag zwischen 0,455 und 0,637. Es zeigt,
wie stark der berechnete Wert bei einer erneuten
Stichprobenziehung nach diesem Verfahren
schwanken kann. Das Intervall berücksichtigt
jedoch nicht alle Unsicherheiten des Bewegungsmodells.
Insbesondere können mehrere virtuelle Kameraeintritte
aus demselben simulierten Bewegungsverlauf stammen.
Diese Abhängigkeit wird beim Ziehen einzelner
Ereigniszeiten nicht berücksichtigt.
Eine weitere Einschränkung betrifft die Datenbasis:
Ein Teil der realen Fotozeitpunkte stammt von
Gerätekennungen ohne bekannte Koordinaten.
Diese Aufnahmen können zeitlich berücksichtigt
werden, für die Simulation stehen dagegen nur
die bekannten Kamerastandorte zur Verfügung.
Die beiden Datengrundlagen sind deshalb räumlich
nicht vollständig identisch.

Auch der angenommene Kameraradius und die
KDE-Bandbreite beeinflussen die Auswertung.
Der Radius legt fest, wann eine virtuelle
Annäherung erfasst wird. Die Bandbreite bestimmt,
wie stark die zeitlichen Verteilungen geglättet
werden. Beide Werte sind Modellparameter
und keine aus japanischen Bärenbewegungen
bestimmten Größen.

Die Auswertung zeigt somit eine teilweise
zeitliche Übereinstimmung zwischen realen
Fotoaufnahmen und simulierten Kameraeintritten.
Sie stellt aber keinen Nachweis dafür dar,
dass die simulierten Bewegungswege das tatsächliche
Verhalten japanischer Bären abbilden.
Da der Vergleich ohne Geländegewichtung
durchgeführt wurde, lässt sich daraus auch
keine Aussage über deren Einfluss auf die
zeitliche Übereinstimmung ableiten.


\section{Bewertung der Anforderungen}
\label{sec:anforderungsbewertung}

Zum Abschluss werden die in
Abschnitt~\ref{sec:anforderungen}
festgelegten Anforderungen mit den
Ergebnissen der Implementierung und
der Funktionsprüfungen abgeglichen.
Dabei wird zwischen umgesetzten Funktionen,
teilweise überprüften Anforderungen
und noch offenen Nachweisen unterschieden.

\subsection*{Funktionale Anforderungen}

Tabelle~\ref{tab:bewertung-funktional}
zeigt die Bewertung der funktionalen
Anforderungen F1 bis F8.

\begin{table}[htbp]
    \centering
    \caption{Bewertung der funktionalen Anforderungen}
    \label{tab:bewertung-funktional}
    \small
    \setlength{\tabcolsep}{4pt}
    \begin{tabular}{L{1.0cm} L{3.8cm} L{2.5cm} L{7.0cm}}
        \toprule
        \textbf{Nr.} &
        \textbf{Anforderung} &
        \textbf{Bewertung} &
        \textbf{Begründung} \\
        \midrule

        F1 &
        Gemeinsame Kartenansicht &
        Erfüllt &
        Die Hauptkarte stellt die verschiedenen
        Datenebenen gemeinsam dar. \\

        \addlinespace
        F2 &
        Datenfilterung &
        Erfüllt &
        GPS-Daten lassen sich nach Bär,
        Zeitraum und Jahresphase filtern. \\

        \addlinespace
        F3 &
        Karteninteraktion &
        Erfüllt &
        Das Anlegen, Speichern und Löschen
        eigener Marker und Polygone
        wurde erfolgreich getestet. \\

        \addlinespace
        F4 &
        Bewegungsanimation &
        Erfüllt &
        Reale und simulierte Bewegungen
        können synchron abgespielt
        und zeitlich gesteuert werden. \\

        \addlinespace
        F5 &
        Zusammenhängende Simulation &
        Erfüllt &
        Das Modell erzeugt zusammenhängende
        simulierte Bewegungswege. \\

        \addlinespace
        F6 &
        Japan-Simulation &
        Erfüllt &
        Bewegungswege können ohne
        vorgegebene GPS-Trajektorie
        erzeugt werden. \\

        \addlinespace
        F7 &
        Statistischer Vergleich &
        Erfüllt &
        Reale Fotozeitpunkte und virtuelle
        Kameraeintritte können statistisch
        gegenübergestellt werden. \\

        \addlinespace
        F8 &
        Dateneinsicht und neue Aufnahmen &
        Teilweise geprüft &
        Daten und gespeicherte
        Bildinformationen sind zugänglich.
        Der automatische Eingang neuer
        Aufnahmen konnte nicht im laufenden
        Betrieb überprüft werden. \\

        \bottomrule
    \end{tabular}
\end{table}

Die funktionalen Anforderungen konnten
damit überwiegend umgesetzt und überprüft
werden. Bei F8 besteht eine Einschränkung,
da die externe Bildextraktion aufgrund
fehlender Zugriffsberechtigungen derzeit
nicht aktiv ist. Die Anzeige bereits
gespeicherter Daten bleibt davon
unabhängig nutzbar.

\subsection*{Nicht-funktionale Anforderungen}

Tabelle~\ref{tab:bewertung-nichtfunktional}
fasst die Bewertung der nicht-funktionalen
Anforderungen N1 bis N5 zusammen.


\begin{table}[htbp]
    \centering
    \caption{Bewertung der nicht-funktionalen Anforderungen}
    \label{tab:bewertung-nichtfunktional}
    \small
    \setlength{\tabcolsep}{4pt}

    \begin{tabular}{L{1.0cm} L{3.8cm} L{2.5cm} L{7.0cm}}
        \toprule
        \textbf{Nr.} &
        \textbf{Anforderung} &
        \textbf{Bewertung} &
        \textbf{Begründung} \\
        \midrule

        N1 &
        Darstellungsleistung &
        Im Funktionstest erfüllt &
        Mehrere Datenebenen konnten gleichzeitig
        dargestellt werden. Die Hauptkarte
        blieb dabei bedienbar. \\

        \addlinespace
        N2 &
        Wartbarkeit &
        Strukturell erfüllt &
        Die Anwendung ist modular aufgebaut.
        Datenverarbeitung, Simulation und
        Kartendarstellung sind voneinander getrennt. \\

        \addlinespace
        N3 &
        Datenzugriff &
        Weitgehend erfüllt &
        MariaDB sowie lokale CSV- und
        SQLite-Daten werden über
        Zugriffsfunktionen in die
        Anwendung eingebunden. \\

        \addlinespace
        N4 &
        Fehlerbehandlung &
        Teilweise überprüft &
        Für ausgewählte Fehlerfälle sind
        Fehlermeldungen und alternative
        Datenzugriffe vorgesehen. \\

        \addlinespace
        N5 &
        Reproduzierbarkeit &
        Für Testfälle erfüllt &
        Simulationen und statistische
        Auswertungen liefern bei gleichen
        Eingabedaten, Parametern und
        Zufallsseeds reproduzierbare
        Ergebnisse. \\

        \bottomrule
    \end{tabular}
\end{table}



Die nicht-funktionalen Anforderungen wurden
bei der Entwicklung berücksichtigt und
in unterschiedlichem Umfang überprüft.
Die gemeinsame Darstellung mehrerer
Datenebenen, der modulare Aufbau und
die reproduzierbare Durchführung von
Simulationen konnten nachvollzogen werden.

Einschränkungen bestehen bei der
Fehlerbehandlung, da nicht alle möglichen
Ausfallsituationen untersucht wurden.
Zudem wurde die Darstellungsleistung
nicht durch systematische Laufzeitmessungen
bewertet. Die grundlegende Bedienbarkeit
der Hauptkarte konnte jedoch im
Funktionstest bestätigt werden.



% ====================================================================

\chapter{Fazit und Ausblick}
\label{chap:fazit}

\section{Fazit}

Ziel dieser Bachelorarbeit war es, die
bestehenden Teilsysteme zur Darstellung
und Verarbeitung von Tiersichtungsdaten
in einer gemeinsamen Anwendung
zusammenzuführen. Darüber hinaus sollten
reale und simulierte Bärenbewegungen
visualisiert und miteinander verglichen
werden können.

Mit der entwickelten MapLibre-basierten
Kartenplattform wurde eine gemeinsame
Darstellung unterschiedlicher Datenbestände
ermöglicht. Das bestehende Bewegungsmodell
wurde in die Anwendung eingebunden und
um die Erzeugung zusammenhängender
Bewegungswege erweitert. Dadurch lassen
sich reale und simulierte Bewegungen
in den Karpaten zeitlich synchron
gegenüberstellen.

Die Übertragung des Bewegungsmodells
auf das japanische Untersuchungsgebiet
zeigte außerdem, dass sich auch ohne
vorgegebene GPS-Trajektorien mögliche
Bewegungswege erzeugen lassen.
Die ergänzende Geländegewichtung
ermöglicht es, räumliche Eigenschaften
bei der Richtungswahl zu berücksichtigen.

Eine weitere Erkenntnis ergibt sich
aus dem statistischen Vergleich mit
den japanischen Kameraaufnahmen.
Die untersuchten Tageszeitverteilungen
zeigten eine teilweise Übereinstimmung.
Damit konnte eine Möglichkeit geschaffen
werden, zwei unterschiedlich aufgebaute
Datenbestände miteinander zu vergleichen.
Eine zuverlässige Aussage über das
tatsächliche Bewegungsverhalten
japanischer Bären lässt sich daraus
jedoch nicht ableiten.

Der Anforderungsabgleich in
Abschnitt~\ref{sec:anforderungsbewertung}
zeigt, dass die wesentlichen funktionalen
Ziele umgesetzt wurden. Einschränkungen
bestehen insbesondere bei der Überprüfung
der Darstellungsleistung, einzelner
Fehlerfälle und des automatischen
Bildabrufs. Diese Punkte begrenzen
die abschließend nachgewiesenen
Eigenschaften der Anwendung,
ohne die umgesetzten Funktionen
grundsätzlich infrage zu stellen.

Insgesamt wurde mit dieser Arbeit
eine technische Grundlage geschaffen,
die zuvor getrennte Daten und
Funktionen miteinander verbindet.
Sie ermöglicht sowohl die gemeinsame
Visualisierung von Beobachtungsdaten
als auch die Untersuchung simulierter
Tierbewegungen. Dabei zeigt die Arbeit,
dass die technische Umsetzbarkeit
einer Simulation von der
wissenschaftlichen Aussagekraft
ihrer Ergebnisse unterschieden
werden muss.






\section{Ausblick}

Die entwickelte Anwendung bietet verschiedene
Möglichkeiten für zukünftige Erweiterungen.

\paragraph{Untersuchung verschiedener Simulationsparameter.}
Zukünftig könnte untersucht werden, wie sich
unterschiedliche Simulationsparameter auf die
statistischen Ergebnisse auswirken. Insbesondere
wäre interessant, ob kleinere Erfassungsradien
der Kameras die Anzahl der simulierten
Kameraereignisse und den Überlappungskoeffizienten
verändern.

\paragraph{Realistischere Modellierung der Kameraerfassung.}
Die bisherige Simulation verwendet kreisförmige
Erfassungsbereiche um die Kamerastandorte.
Zukünftig könnte stattdessen berücksichtigt
werden, in welche Richtung eine Kamera blickt
und welchen Bereich vor ihr sie tatsächlich
erfasst. Dafür wären Informationen über
Blickrichtung, Erfassungswinkel und Reichweite
der eingesetzten Kameras erforderlich.
Auch mögliche Unterschiede zwischen Tag und
Nacht könnten untersucht werden.

\paragraph{Optimierung der 3D-Kartendarstellung.}
Eine weitere Möglichkeit besteht darin,
die Darstellungsleistung der dreidimensionalen
Geländeansicht zu verbessern. Insbesondere
könnten die Ladezeiten und die Verarbeitung
der Geländedaten genauer untersucht werden.

\paragraph{Auswertung einzelner Kamerastandorte.}
Wenn zukünftig mehr reale Bärenaufnahmen
vorliegen, könnten die zeitlichen
Aktivitätsverteilungen auch für einzelne
Kamerastandorte untersucht werden.
Dadurch wäre ein Vergleich der
Aktivitätsmuster zwischen verschiedenen
Standorten möglich.

Diese Ansätze bieten Möglichkeiten,
die Simulation und die statistische
Auswertung in zukünftigen Arbeiten
weiterzuentwickeln.



% ====================================================================
\clearpage
\printbibliography[
    heading=bibintoc,
    title={Literaturverzeichnis}
]

\clearpage
\thispagestyle{empty}


\begin{singlespace}

\section*{Erklärung zur Verwendung KI-gestützter Hilfsmittel}

Bei der Erstellung dieser Bachelorarbeit wurden
die KI-Werkzeuge unterstützend eingesetzt.

Die Nutzung umfasste insbesondere:

\begin{itemize}
    \item Sprachliche Korrekturen und Verbesserung
    von Formulierungen.

    \item Unterstützung bei der Fehlersuche,
    Optimierung und Bereinigung von Programmcode.

    \item Unterstützung bei der Literaturrecherche
    und Überprüfung bibliografischer Angaben.

    \item Erstellung synthetischer Testbilder
    zur Überprüfung der automatischen
    Bildextraktion aus einem E-Mail-Postfach.


    \item Übersetzung von Texten und Bezeichnungen
    der Anwendung ins Englische.

    \item Unterstützung bei der Strukturierung
    und Überarbeitung einzelner Textabschnitte.
\end{itemize}

Die KI-Werkzeuge dienten als Hilfsmittel bei der
Bearbeitung der Aufgaben. Die fachliche
Verantwortung für die Inhalte, die Entwicklung
der Anwendung und die Bewertung der Ergebnisse
liegt beim Verfasser.

\vspace{2cm}

Regensburg, \today

\vspace{1.5cm}

\rule{7cm}{0.4pt}

Yassir Sbai

\end{singlespace}



\end{document}